\exercisesection{}

\begin{enumerate}
\def\labelenumi{\arabic{enumi})}
\tightlist
\item
  设 \(\mathbf{E}\) 为一个集合，\(\Phi\) 为 \(\mathbf{E}\) 的子集所成的一个集环，\(m\) 为定义在 \(\Phi\) 上的一个可加集函数。用 \(\mathcal{E}(\Phi)\) 表示实 \(\Phi\) -阶梯函数在 \(\mathbf{E}\) 上所成的向量空间，并用 \(J\) 表示 \(\mathcal{E}(\Phi)\) 上满足 \(J(\varphi_{\mathbf{A}}) = m(\mathbf{A})\)（对所有 \(A \in \Phi\)）的线性形式。
\end{enumerate}

\begin{enumerate}
\def\labelenumi{\alph{enumi})}
\tightlist
\item
  对每个 \(\mathbf{A} \in \Phi\)，置
\end{enumerate}

\[
m _ {+} (\mathrm{A}) = \sup _ {\mathrm{B} \in \Phi , \mathrm{B} \subset \mathrm{A}} m (\mathrm{B}) \quad \text { and } \quad m _ {-} (\mathrm{A}) = \sup _ {\mathrm{B} \in \Phi , \mathrm{B} \subset \mathrm{A}} (- m (\mathrm{B})).
\]

证明 \(|m|(\mathbf{A}) = m_{+}(\mathbf{A}) + m_{-}(\mathbf{A})\) 是形如 \(\sum_{i=1}^{p}|m(\mathbf{A}_i)|\) 的数所成集合的上确界（对所有有限划分 \((\mathbf{A}_i)_{1\leqslant i\leqslant p}\)（把 A 划分成属于 \(\Phi\) 的集合））。

\begin{enumerate}
\def\labelenumi{\alph{enumi})}
\setcounter{enumi}{1}
\item
  称 \(m\) 是相对有界的，如果形如 \(m(\mathbf{B})\)（其中 \(\mathbf{B} \in \Phi\)，\(\mathbf{B} \subset \mathbf{A}\)）的数所成的集合对任何集合 \(\mathbf{A} \in \Phi\) 都是有界的。证明下列条件等价：\(\alpha)\) \(m\) 相对有界；\(\beta)\) 线性形式 \(J\) 在 Riesz 空间 \(\mathcal{E}(\Phi)\) 上相对有界；\(\gamma)\) \(m\) 是 \(\Phi\) 上两个正可加集函数之差；\(\delta)\) \(|m|(\mathbf{A})\) 对所有 \(\mathbf{A} \in \Phi\) 有限（按 \(\alpha) \Rightarrow \beta) \Rightarrow \gamma) \Rightarrow \delta) \Rightarrow \alpha)\) 循环论证）。
\item
  假设 \(m\) 相对有界。证明集函数 \(|m|\)、\(m_{+}\) 与 \(m_{-}\) 都是可加的，且 \(m = m_{+} - m_{-}\)（利用 Ch. II, §2, No.~2 的 Th. 1）。此外，若 \(m'\) 与 \(m''\) 是 \(\Phi\) 上满足 \(m = m' - m''\) 的可加正集函数，则 \(m' \geqslant m_{+}\) 且 \(m'' \geqslant m_{-}\)。
\end{enumerate}

\(d)\) 称 \(m\) 是可数可加的，如果 \(m(A) = \sum_{n\geqslant 1}m(A_n)\) 对每个可数

划分 \((\mathbf{A}_n)_{n\geqslant 1}\)（把集合 \(\mathbf{A}\in \Phi\) 划分成属于 \(\Phi\) 的集合）都成立。证明这一条件意味着：\(\lim_{n\to \infty}m(\mathbf{A}_n) = 0\) 对每个递减序列 \((\mathbf{A}_n)_{n\geqslant 1}\)（由 \(\Phi\) 的元素组成、满足 \(\bigcap_{n\geqslant 1}\mathbf{A}_n = \emptyset\)）成立。若 \(m\) 可数可加，证明 \(|m|\) 也可数可加。

\begin{enumerate}
\def\labelenumi{\arabic{enumi})}
\setcounter{enumi}{1}
\tightlist
\item
  设 E 为一个集合，\(\Phi\) 为 E 的子集所成的一个集环，V 为一个全格序的 Riesz 空间。称映射 \(m\)（\(\Phi\) 到 V 中的）为正的，如果有 \(m(A) \geqslant 0\)（对所有 \(A \in \Phi\)）；称其为可加的，如果有 \(m(A \cup B) = m(A) + m(B)\)（当集合 \(A \in \Phi\) 与 \(B \in \Phi\) 不相交时）；称其为相对有界的，如果元素 \(m(B)\)（其中 \(B \in \Phi\)，\(B \subset A\)）所成的集合在 V 中有上界（对任何 \(A \in \Phi\)）。
\end{enumerate}

\begin{enumerate}
\def\labelenumi{\alph{enumi})}
\tightlist
\item
  设 \(m\) 是 \(\Phi\) 到 \(V\) 中的一个相对有界映射。对每个 \(A\)，按 Exer. 1 a) 的方式定义元素 \(m_{+}(A), m_{-}(A)\) 和 \(|m|(A)\)（均属于 \(V\)）。证明 \(|m|(A)\)
\end{enumerate}

是 V 中形如 \(\sum_{i=1}^{p}|m(A_i)|\) 的元素所成集合的上确界（对所有有限划分 \((\mathbf{A}_i)_{1\leqslant i\leqslant p}\)（把 A 划分成属于 \(\Phi\) 的集合））。由此推出映射 \(|m|\)（\(\Phi\) 到 V 中的）是可加的，继而推出 \(m\) 是两个可加正映射 \(m_{+}\) 与 \(m_{-}\)（均由 \(\Phi\) 到 V 中）之差。

\begin{enumerate}
\def\labelenumi{\alph{enumi})}
\setcounter{enumi}{1}
\tightlist
\item
  设 \(m'\) 与 \(m''\) 是 \(\Phi\) 到 \(V\) 中的两个可加正映射，且 \(m = m' - m''\)。证明 \(m\) 相对有界，且 \(m' \geqslant m_+\)，\(m'' \geqslant m_-\)。c) 推广 Exer. 1 d)。
\end{enumerate}

\begin{enumerate}
\def\labelenumi{\arabic{enumi})}
\setcounter{enumi}{2}
\tightlist
\item
  设 \(\mathbf{E}\) 为一个集合，\(\mathfrak{T}\) 为 \(\mathbf{E}\) 的子集所成的一个 σ-集环。用 \(\mathcal{M}\) 表示映射 \(f\)（\(\mathbf{E}\) 到 \(\overline{\mathbf{R}}_+\) 中的）所成的集合，这些映射使得 \(x \in \mathbf{E}\) 中使 \(f(x) \geqslant c\) 的 x 所成的集合属于 \(\mathfrak{T}\)（对每个 \(c \in \overline{\mathbf{R}}_+\)）。
\end{enumerate}

\begin{enumerate}
\def\labelenumi{\alph{enumi})}
\item
  证明 \(\mathcal{M}\) 中元素组成的每个序列的上极限和下极限都属于 \(\mathcal{M}\)（先研究递增或递减的序列）。
\item
  证明 \(\mathfrak{T}\) 是 \(\mathbf{E}\) 中这样的子集所成的集合：其特征函数属于 \(\mathcal{M}\)。
\item
  设 \(f \in \mathcal{M}\)。对每个 \(B\)（\(\overline{\mathbf{R}}_+\) 的 Borel 子集），集合 \(f^{-1}(B)\) 属于 \(\mathfrak{T}\)（子集 \(B\)（属于 \(\overline{\mathbf{R}}_+\)）中使 \(f^{-1}(B) \in \mathfrak{T}\) 的那些 B 所成的集合是一个 σ-集环，且包含区间 \([c, +\infty]\)）。
\item
  设 \(f_1, \ldots, f_n\) 属于 \(\mathcal{M}\)，\(\varphi\) 是 \((\overline{\mathbf{R}}_+)^n\) 到 \(\overline{\mathbf{R}}_+\) 中的一个 Borel 映射；证明映射 \(x \mapsto \varphi(f_1(x), \ldots, f_n(x))\)（E 到 \(\overline{\mathbf{R}}_+\) 中的）属于 \(\mathcal{M}\)。由此推出 \(\mathcal{M}\) 包含由 \(\mathcal{M}\) 的元素组成的每个级数的和。
\item
  证明 \(\mathcal{M}\) 是有限正 \(\mathfrak{T}\) -阶梯函数组成的递增序列的极限所成的集合。
\end{enumerate}

\begin{enumerate}
\def\labelenumi{\arabic{enumi})}
\setcounter{enumi}{3}
\tightlist
\item
  记号与假设同上一习题。所谓 \((\mathbf{E},\mathfrak{T})\) 上的抽象测度，是指每个映射 \(m\)（从 \(\mathfrak{T}\) 到 \(\overline{\mathbf{R}}_+\) 中，且具有下列性质）：对每个序列 \((\mathrm{A}_n)_{n\in \mathbf{N}}\)（由属于 \(\mathfrak{T}\) 的集合组成，且两两不相交），有 \(m\big(\bigcup_{n\in \mathbf{N}}\mathrm{A}_n\big) =\)
\end{enumerate}

\(\sum_{n\in \mathbf{N}}m(\mathrm{A}_n)\)。所谓 \((\operatorname {E},\mathfrak{T})\) 上的积分，是指 \(\mathcal{M}\) 到 \(\overline{\mathbf{R}}_+\) 中且满足下列条件的每个映射 J：

  下列条件：\(\alpha)\) \(J(\lambda \cdot f) = \lambda \cdot J(f)\) 对 \(\lambda \in \overline{\mathbf{R}}_{+}\) 与 \(f \in \mathcal{M}\) 成立（按通常的约定 \(0 \cdot (+\infty) = 0\)）；\(\beta)\) \(J\big(\sum_{n \in \mathbf{N}} f_n\big) = \sum_{n \in \mathbf{N}} J(f_n)\) 对每个序列 \((f_n)_{n \in \mathbf{N}}\)（由 \(\mathcal{M}\) 的元素组成）成立。


\begin{enumerate}
\def\labelenumi{\alph{enumi})}
\item
  设 \(J\) 是 \((\mathbf{E},\mathfrak{T})\) 上的一个积分；对每个子集 \(A \in \mathfrak{T}\)，置 \(m_{J}(A) = J(\varphi_A)\)；证明 \(m_J\) 是一个抽象测度，且映射 \(J \mapsto m_J\) 是 \((\mathbf{E},\mathfrak{T})\) 上积分所成的集合到 \((\mathbf{E},\mathfrak{T})\) 上抽象测度所成的集合上的一个双射（若 \(m\) 是抽象测度，先定义 \(J(f)\)（按线性，对有限正 \(\mathfrak{T}\) -阶梯函数），并利用 Exer. 3 e)）。若 \(m\) 是抽象测度，我们将用 \(m^*\) 表示对应的积分。
\item
  设 \(m\) 是 (E, T) 上的一个抽象测度。对 E 上的每个正函数 \(f\)（无论有限与否），置 \(m^{*}(f) = \inf_{g \in \mathcal{M}, g \geqslant f} m^{*}(g)\)。证明 \(m^{*}\) 是 E 上的一个负担（模仿 Ch. IV, §1, No.~3 的 Th. 3 的证明）。对 E 的每个子集 A，置 \(m^{*}(A) = m^{*}(\varphi_{A})\)。
\end{enumerate}

\begin{enumerate}
\def\labelenumi{\arabic{enumi})}
\setcounter{enumi}{4}
\tightlist
\item
  设 \(\mathbf{E}\) 为一个集合，\(\mathfrak{T}\) 为 \(\mathbf{E}\) 的子集所成的一个 σ-集环，\(m\) 为 \((\mathbf{E},\mathfrak{T})\) 上的一个抽象测度，\(p\geqslant 1\) 为一个有限实数。对每个数值函数 \(f\)（\(\mathbf{E}\) 上的），置 \(\mathrm{N}_p(f) = m^* (|f|^p)^{1 / p}\)，并用 \(\mathcal{F}^p\) 表示函数 \(f\)（使 \(\mathrm{N}_p(f)\) 为有限）所成的集合（cf.~§1, Exer. 4）。给向量空间 \(\mathcal{F}^p\) 赋以半范数 \(\mathrm{N}_p\)，对它而言该空间是完备的。
\end{enumerate}

\begin{enumerate}
\def\labelenumi{\alph{enumi})}
\item
  设 \(\mathcal{N}\) 为使 \(m^{*}(|f|)=0\) 的函数 f 所成的集合，\(\mathcal{V}\) 为由属于 \(\mathcal{M}\) 的有限函数生成的向量空间，\(\mathcal{L}^{p}\) 为 \(\mathcal{F}^{p}\) 中包含集合 \(A\in \mathfrak{T}\)（使 \(m^{*}(A)\) 为有限）的特征函数的最小闭线性子空间。证明 \(\mathcal{L}^{p}=\mathcal{N}+(\mathcal{V}\cap\mathcal{F}^{p})\)。
\item
  把 Lebesgue 定理（Ch. IV, §3, No.~7, Th. 6）推广到目前的情形，并特别考察 \(p = 1\) 的情形（此时应定义函数 \(f \in \mathcal{L}^1\) 的积分）。
\item
  称 \(\mathbf{A}\)（\(\mathbf{E}\) 的一个子集）是 \(m\) -可积的，如果 \(\varphi_{\mathbf{A}} \in \mathcal{L}^1\)。证明此时存在集合 \(\mathbf{A}' \in \mathfrak{T}\)，使得 \(m^{*}(\mathbf{A} \cap \mathbb{C}\mathbf{A}') = m^{*}(\mathbf{A}' \cap \mathbb{C}\mathbf{A}) = 0\)。在 \(m(\mathbf{E})\) 有限的情形，逆命题成立。
\end{enumerate}

¶ 6) 设 E 为一个集合，T 为 E 的子集所成的一个 σ-集环，m 为 (E, T) 上的一个抽象测度。假设 \(m(E)\) 有限，并且对每个 \(B \in T\)（满足 \(m(B) > 0\)），存在集合 \(A \in T\)，使得 \(A \subset B\) 且 \(0 < m(A) < m(B)\)。则对每个 \(B \in T\) 与每个满足 \(0 \leqslant t \leqslant m(B)\) 的数 t，存在集合 \(A \in T\)，使得 \(m(A) = t\) 且 \(A \subset B\)。

¶ 7) 设 T 为一个拓扑空间，\(\Phi\) 为 T 的子集所成的一个集环，\(m\) 为 \(\Phi\) 到 R 中的一个相对有界可加映射（cf.~Exer. 1）。假设存在 T 的开覆盖 \(\mathfrak{U}\)，使得 \(\Phi\) 由包含在 \(\mathfrak{U}\) 中有限多个元素的并内的 T 的 Borel 子集组成。再假设对每个 \(A \in \Phi\) 与每个 \(\varepsilon > 0\)，存在 T 的紧子集 K 与 T 的开子集 U，使得 \(K \subset A \subset U\)，\(U \in \Phi\) 且 \(|m|(U - K) < \varepsilon\)。证明 T 上存在测度 \(\mu\)（不必为正），使得 \(m(A) = \mu^{\bullet}(A)\)（对所有 \(A \in \Phi\)），并且这样的测度 \(\mu\) 是唯一的（通过引入 Exer. 1 c) 的分解 \(m = m_{+} - m_{-}\)，化为 \(m\) 取正值的情形；先处理 \(T \in \Phi\) 的情形，并注意此时 \(m\) 是内正则的；最后通过把测度粘合起来处理一般情形）。

\begin{enumerate}
\def\labelenumi{\arabic{enumi})}
\setcounter{enumi}{7}
\tightlist
\item
  设 T 为一个拓扑空间，m 为 \(\mathfrak{B}(\mathrm{T})\) 到 R 中的一个有界、可数可加的映射（cf.~Exer. 1d)）。用 T 表示具有下列性质的 T 的 Borel 子集 A 所成的集合：对每个 \(\varepsilon > 0\)，存在 T 中的闭集 F 与开集 U，使得 \(F \subset A \subset U\)，并且对 U - F 的每个 Borel 子集 B 有 \(|m(B)| < \varepsilon\)（换句话说，\(|m|(\mathrm{U} - \mathrm{F}) < \varepsilon\)）。证明 T 是 T 的子集所成的一个 σ-集环（先注意 \(T \in T\)，且 T 是一个集环；然后只须证明：若 \((\mathrm{A}_{n})_{n \in \mathbb{N}}\) 是 T 中两两不相交的元素组成的序列，则集合 \(A = \bigcup_{n \in N} A_{n}\) 属于 T；为此，选取闭集 \(F_{n}\) 与开集 \(U_{n}\)，使得 \(F_{n} \subset A_{n} \subset U_{n}\) 且 \(|m|(\mathrm{U}_{n} - \mathrm{F}_{n}) < \varepsilon / 2^{n}\)，然后对充分大的 p 置 \(F = F_{0} \cup \cdots \cup F_{p}\)，并置 \(U = \bigcup_{n \in \mathbf{N}} U_{n}\)。
\end{enumerate}


\begin{enumerate}
\def\labelenumi{\arabic{enumi})}
\setcounter{enumi}{8}
\tightlist
\item
  设 \(\mathbf{X}\) 为一个集合；称 \(\mathbf{X}\) 上的每个可数可加映射（\(\mathfrak{P}(\mathbf{X})\) 到 \(\mathbf{R}_+\) 中的）为负担。若负担 \(m\)（\(\mathbf{X}\) 上的）满足 \(m(\{x\}) = 0\)（对所有 \(x \in \mathbf{X}\)），则称它是弥散的；若存在正函数 \(f\)（在 \(\mathbf{X}\) 上），使得 \(m(\mathbf{A}) = \sum_{x \in \mathbf{A}} f(x)\) 对每个
\end{enumerate}

X 的子集 A 都成立，则称其为原子的。

\begin{enumerate}
\def\labelenumi{\alph{enumi})}
\item
  证明 X 上的每个负担都可以以唯一的方式分解为一个弥散负担与一个原子负担之和。
\item
  设 \(X'\) 为 \(X\) 的一个子集，\(m'\) 为 \(X'\) 上的一个负担；对每个子集 \(A\)（\(X\) 的），置 \(m(A) = m'(A \cap X')\)。证明 \(m\) 是 \(X\) 上的一个负担，并且它是弥散的（相应地，原子的）当且仅当 \(m'\) 具有同样的性质。
\item
  设 \(m\) 为 \(X\) 上的一个负担，\((X_i)_{i \in I}\) 为 \(X\) 的两两不相交的子集组成的一个族。若 \(I\) 上的每个负担都是原子的，则 \(m\big(\bigcup_{i \in I} X_i\big) = \sum_{i \in I} m(X_i)\)。
\end{enumerate}

\begin{enumerate}
\def\labelenumi{\arabic{enumi})}
\setcounter{enumi}{9}
\tightlist
\item
  设 \(\mathbf{X}\) 为一个不可数的无限集，带有一个良序关系，记作 \(x \leqslant y\)。对每个 \(x \in \mathbf{X}\)，用 \(\mathrm{I}(x)\) 表示 \(y \in \mathbf{X}\) 中满足 \(y < x\) 的 y 所成的集合。我们作下列假设：\(\alpha\) ) 存在一个最大元素 \(a\)（在 \(\mathbf{X}\) 中）；\(\beta\) ) 严格小于 \(\operatorname{Card}(\mathbf{X})\) 的基数所成的集合有一个最大元素 \(\mathfrak{c}\)；\(\gamma\) ) 对每个 \(x\)（在 \(\mathrm{I}(a)\) 中），有 \(\operatorname{Card}(\mathrm{I}(x)) \leqslant \mathfrak{c}\)，且集合 \(\mathbf{Y}\)（由 \(x \in \mathbf{X}\) 中使 \(\operatorname{Card}(\mathrm{I}(x)) = \mathfrak{c}\) 的 x 组成）非空；
\end{enumerate}

\(\delta)\) 基数 \(\leqslant \mathfrak{c}\) 的集合上的每个负担都是原子的。证明 X 上的每个负担都是原子的。可以按如下方式用反证法论证：设 \(m\) 为 X 上的一个非零弥散负担。用 \(b\) 表示 Y 的最小元素，并对每个 \(x \in I(a)\)，令 \(f_x\) 为 \(I(x)\) 到 \(I(b)\) 中的一个单射。对每对 \((x,y) \in I(a) \times I(b)\)，用 \(A_{x,y}\) 表示由元素 \(z \in X\)（满足 \(x < z < a\) 且 \(f_z(x) = y\)）所成的集合；对每个 \(x \in I(a)\)，用 \(M_x\) 表示由 \(y \in I(b)\)（满足 \(m(A_{x,y}) > 0\)）所成的集合，并对每个 \(y \in I(b)\)，用 \(N_y\) 表示由 \(x \in I(a)\)（满足 \(m(A_{x,y}) > 0\)）所成的集合。利用 Exer. 9 c)，证明 \(m(X) = \sum_{y \in I(b)} m(A_{x,y})\)；由此推出集合 \(M_x\) 对每个 \(x \in I(a)\) 都是可数的且非空；再证明每个集合 \(N_y\) 都是可数的，并由此得出矛盾 \(\text{Card}(X) = c\)。

\begin{enumerate}
\def\labelenumi{\arabic{enumi})}
\setcounter{enumi}{10}
\tightlist
\item
  称基数 \(\mathfrak{c}\) 为 Ulam 基数（或 Ulam 型的），如果每个具有基数 \(\mathfrak{c}\) 的集合上的负担都是原子的。证明每个可数基数都是 Ulam 型的；若 \(\mathfrak{c}\) 是 Ulam 基数，则每个基数 \(\mathfrak{c}' \leqslant \mathfrak{c}\) 也是 Ulam 基数；并且若 \((\mathfrak{c}_i)_{i \in I}\) 是一族 Ulam 基数，且 \(\operatorname{Card}(I)\) 是 Ulam 型的，则基数 \(\sum_{i \in I} \mathfrak{c}_i\) 是 Ulam 型的。最后，
\end{enumerate}

最小不可数基数是 Ulam 型的（cf.~Exer. 10）。

\begin{enumerate}
\def\labelenumi{\arabic{enumi})}
\setcounter{enumi}{11}
\tightlist
\item
  设 X 为一个集合，U 为 \(\mathfrak{P}(X)\) 的一个子集，m 为 U 的特征函数。为使 m 是一个负担，必要且充分的是 U 是一个超滤子，并且 U 的每个可数子集的交属于 U。在这些条件下，称 U 为 X 上的一个 Ulam 超滤子。假设情形如此，且 \((\mathrm{H}_{i})_{i\in\mathrm{I}}\) 是 U 中元素组成的一个族，其中 \(\operatorname{Card}(\mathrm{I})\) 是 Ulam 基数；证明 \(\bigcap_{i\in I} H_{i}\) 属于
\end{enumerate}

\(\mathfrak{U}\)。

\begin{enumerate}
\def\labelenumi{\arabic{enumi})}
\setcounter{enumi}{12}
\tightlist
\item
  在本习题中，我们承认连续统假设，也就是说，我们把公理 \(\langle\mathrm{Card}(\mathbf{R})\) 是最小不可数基数 \(\rangle\) 附加到集合论的公理中。\(^{(1)}\)
\end{enumerate}

\begin{enumerate}
\def\labelenumi{\alph{enumi})}
\item
  \(\mathbf{R}\) 上的每个负担都是原子的（cf.~Exer. 11）。
\item
  设 X 为一个集合，m 为 X 上满足下列性质的负担：对每个满足 \(m(A) > 0\) 的 X 的子集 A，存在 A 的子集 B，使得 \(0 < m(B) < m(A)\)；则 m 为零（对每个整数 \(n \geqslant 1\)，存在 X 的一个有限划分 \((\mathrm{X}_{n,i})_{i \in \mathrm{I}_n}\)，使得 \(m(\mathrm{X}_{n,i}) \leqslant 1/n\)（对所有 \(i \in I_n\)）；置 \(I = \prod_{n \geqslant 1} I_n\) 与
\end{enumerate}

\(\mathrm{X}_{\alpha} = \bigcap_{n\geqslant 1}\mathrm{X}_{n,\alpha_n}\)（对 I 中的每个 \(\alpha = (\alpha_{n})_{n\geqslant 1}\)）；族 \((\mathrm{X}_i)_{i\in \mathrm{I}}\) 是 \(\mathbf{X}\) 的一个划分，

\(m(X_{i}) = 0\)（对每个 \(i \in I\)），并且 \(I\) 上的每个负担由 \(a\) ) 知都是原子的；借助 Exer. 9 c) 完成证明。

\begin{enumerate}
\def\labelenumi{\alph{enumi})}
\setcounter{enumi}{2}
\item
  设 X 是一个不存在非平凡 Ulam 超滤子的集合；则 X 上的每个负担都是原子的（否则，由 b)，X 上将存在一个弥散负担 m 与 X 的子集 A，使得 \(m(A) > 0\)，且对 A 的每个子集 B，\(m(B) = 0\) 或 \(m(A - B) = 0\)；由满足 \(m(A) = m(A \cap B)\) 的 X 的子集 B 组成的集合 U 将是 X 上的一个非平凡 Ulam 超滤子）。
\item
  证明 \(2^{\mathfrak{c}}\) 对每个 Ulam 基数 \(\mathfrak{c}\) 都是 Ulam 基数（由 \(c\) )，只须证明：若 \(C\) 是具有基数 \(\mathfrak{c}\) 的集合，则每个 Ulam 超滤子 \(\mathfrak{U}\)（在 \(X = \{0,1\}^{C}\) 上）都是平凡的；给 \(C\) 装备一个良序结构，并用超限归纳定义一个族 \(s = (s_{\alpha})_{\alpha \in C}\)（由 \(\{0,1\}\) 中元素组成），使得由 \(x = (x_{\alpha})_{\alpha \in C}\)（\(X\) 中满足 \(x_{\alpha}=s_{\alpha}\)（对所有 \(\alpha\leqslant\beta\)）者）组成的集合属于 \(\mathfrak{U}\)，对每个 \(\beta\in\mathrm{C}\)（cf.~Exer. 12）；于是 \(s\) 属于 \(\mathfrak{U}\) 的每个元素）。\(^{(2)}\)
\end{enumerate}

\begin{enumerate}
\def\labelenumi{\arabic{enumi})}
\setcounter{enumi}{13}
\item
  设 \((\mathbf{T}_i)_{i\in I}\) 是一族 Radon 拓扑空间，\(\mathbf{T}\) 是该族的和空间。若 \(\operatorname{Card}(\mathbf{I})\) 是 Ulam 基数，则拓扑空间 \(\mathbf{T}\) 是 Radon 空间（证明：对每个正的、有界的可数可加函数 \(m\)（在 \(\mathbf{T}\) 的 Borel σ-集环上），存在可数子集 \(J\)（\(I\) 的），使得 \(m(\bigcup_{i \in I - J} \mathbf{T}_i) = 0\)）。
\item
  设 T 是一个离散空间，其基数是最小不可数基数。证明 T 是一个 Radon 局部紧空间，其拓扑没有可数基（cf.~Exer. 11）。由此推出存在不可度量化的紧 Radon 空间。
\item
  设 \(I\) 是一个可数集，\((\mathbf{T}_i)_{i \in I}\) 是一族 Radon 空间。假设这些空间 \(\mathbf{T}_i\) 中每一个的紧子集都是可度量化的。证明乘积空间 \(\mathbf{T} = \prod_{i \in I} \mathbf{T}_i\) 是 Radon 空间。推广到逆极限的情形。
\item
  设 \(\mathbf{T}\) 是一个拓扑空间。
\end{enumerate}

\begin{enumerate}
\def\labelenumi{\alph{enumi})}
\tightlist
\item
  设 \(m\) 是 \(\mathfrak{B}(\mathbf{T})\) 到 \(\mathbf{R}_+\) 中的一个可数可加映射。假设存在一个 Borel 子集的序列 \((\mathbf{T}_n)_{n\in \mathbf{N}}\)（\(\mathbf{T}\) 的），使得 \(\mathbf{T} = \bigcup_{n\in \mathbf{N}}\mathbf{T}_n\)，\(m(\mathbf{T}_0) = 0\)
\end{enumerate}

并且子空间 \(\mathbf{T}_n\)（\(\mathbf{T}\) 的）对每个 \(n \geqslant 1\) 都是 Radon 的。证明存在一个有界正测度 \(\mu\)（在 \(\mathbf{T}\) 上），使得 \(m(\mathbf{A}) = \mu^{\bullet}(\mathbf{A})\)（对每个 Borel 子集 \(\mathbf{A}\)（\(\mathbf{T}\) 的））（可以借助 No.~3 的 Prop. 2 的推论，化为这些 \(\mathbf{T}_n\) 两两不相交的情形）。

\begin{enumerate}
\def\labelenumi{\alph{enumi})}
\setcounter{enumi}{1}
\item
  把上一结果推广到如下情形：集合 \(T_{n}\)（对 \(n \geqslant 1\)）不再假设是 Borel 的，而只是 T 中普遍可测的（按 No.~3 的 Prop. 2 的方式论证）。
\item
  每个是一个可度量化紧子空间序列之并的拓扑空间都是 Radon 的。
\end{enumerate}

\begin{enumerate}
\def\labelenumi{\arabic{enumi})}
\setcounter{enumi}{17}
\tightlist
\item
  设 \(\mathbf{T}_1\) 与 \(\mathbf{T}_2\) 是两个拓扑空间，\(f\) 是 \(\mathbf{T}_1\) 到 \(\mathbf{T}_2\) 上的一个连续双射映射。假设对每个有界可数可加函数 \(m_1\)（在 Borel σ-集环 \(\mathfrak{B}(\mathbf{T}_1)\) 上），存在一个由紧子集 \(\mathrm{K}_n\)（\(\mathbf{T}_1\) 的）组成的序列，使得 \(m_1(\mathbf{T}_1 - \bigcup_{n}\mathbf{K}_n) = 0\)。设 \(m_2\) 是 \(\mathfrak{B}(\mathbf{T}_2)\) 上的一个有界可数可加函数。
\end{enumerate}

为使存在一个有界可数可加函数 \(m_{1}\)（在 \(\mathfrak{B}(\mathrm{T}_{1})\) 上），使得 \(m_{2}(\mathrm{A}) = m_{1}(f^{-1}(\mathrm{A}))\)（对所有 \(\mathrm{A} \in \mathfrak{B}(\mathrm{T}_{2})\)），必要且充分的是下列条件成立：存在一个由紧子集 \(K_{n}\)（\(T_{1}\) 的）组成的序列，使得 \(m_{2}(\mathrm{T}_{2} - \bigcup_{n} f(\mathrm{K}_{n})) = 0\)。

\exercisesection{}

\begin{enumerate}
\def\labelenumi{\arabic{enumi})}
\tightlist
\item
  设 \(\mathcal{T} = (\mathrm{K}_i, p_{ij})\) 是以集合 I 为指标的紧空间的一个逆系统，设 K 是一个紧空间，并设 \((p_i)_{i \in I}\) 是连续映射 \(p_i: \mathrm{K} \to \mathrm{K}_i\) 组成的一个相容且分离的族。
\end{enumerate}

\begin{enumerate}
\def\labelenumi{\alph{enumi})}
\item
  设 A 是 K 上形如 \(f_{i} \circ p_{i}\) 的连续函数所成的集合，其中 i 取遍 I，\(f_{i}\) 取遍 \(K_{i}\) 上连续实函数所成的集合。证明 A 是 \(\mathcal{C}(K)\) 的一个稠密线性子空间。
\item
  设 \((\mu_i)_{i \in I}\) 是 \(\mathcal{T}\) 上正测度组成的一个逆子系统。假设这些 \(p_i\) 都是满射。对每个 \(f \in A\)，用 \(I_f\) 表示这样的 \(i \in I\) 组成的集合：\(f\) 形如 \(f_i \circ p_i\)，其中 \(f_i \in \mathcal{C}(K_i)\)（必然唯一）。证明存在唯一的测度 \(\pi\)（在 \(K\) 上），使得 \(\pi(f) = \inf_{i \in I_f} \mu_i(f_i)\)（对所有 \(f \in A\)）。证明 \(\pi\) 是
\end{enumerate}

K 上正测度 \(\mu\)（满足 \(p_i(\mu) \leqslant \mu_i\)，对所有 \(i \in I\)）中最大的一个。

¶ 2) 假设与记号同 No.~2 的 Th. 1，我们打算对它给出一个新的证明。

\begin{enumerate}
\def\labelenumi{\alph{enumi})}
\item
  设 K 是 T 的一个紧子集；对每个 \(i \in I\)，置 \(\mathrm{K}_{i} = p_{i}(\mathrm{T})\)，并用 \(q_{i}\) 表示 K 到 \(K_{i}\) 上的、在 K 上与 \(p_{i}\) 重合的映射；对 \(i \leqslant j\)，用 \(q_{ij}\) 表示 \(K_{j}\) 到 \(K_{i}\) 中的、在 \(K_{j}\) 上与 \(p_{ij}\) 重合的映射。由上一习题推出，存在 K 上最大的正测度 \(\pi^{K}\)，使得 \(q_{i}(\pi^{K}) = (\mu_{i})_{\mathbf{K}_{i}}\)（对所有 \(i \in I\)）。
\item
  若 K 与 L 是 T 的紧子集且 \(K \subset L\)，证明 \((\pi^{\mathrm{L}})_{\mathrm{K}} \geqslant \pi^{\mathrm{K}}\)；由此推出，形如 \(i^{\mathrm{K}}(\pi^{\mathrm{K}})\) 的测度所成的集合（其中 K 取遍 T 的紧子集所成的集合，而 \(i^{K}\) 是 K 到 T 中的典范单射）有一个上确界 \(\mu\)（在 \(\mathcal{M}_{+}(\mathrm{T})\) 中）。
\item
  证明 \(p_i(\mu) = \mu_i\) 对每个 \(i \in I\) 成立（只须注意 \(p_i(\mu) \leqslant \mu_i\)，且由假设 (P)，测度 \(p_i(\mu)\) 与 \(\mu_i\) 有相同的总质量）。
\end{enumerate}

\exercisesection{}

\begin{enumerate}
\def\labelenumi{\arabic{enumi})}
\tightlist
\item
  设 \(\mathbf{T}\) 为一个拓扑空间，\(\mu\) 为 \(\mathbf{T}\) 上的一个正测度。用 \(\mathfrak{S}_{\mu}\) 表示 \(\mathbf{T}\) 中边界为 \(\mu\) -可忽略的子集所成的集合。
\end{enumerate}

\begin{enumerate}
\def\labelenumi{\alph{enumi})}
\item
  证明 \(\mathfrak{S}_{\mu}\) 是一个集环。
\item
  若 T 是完全正则的，则 T 的每个点都有一个包含于 \(S_{\mu}\) 的基本邻域系（注意：若 f 是 T 上的连续函数且在一个 \(\mu\) -可积集合之外为零，则使 \(f(r)\) 不是 \(\mu\) -可忽略的实数 r 所成的集合是可数的）。
\end{enumerate}

\begin{enumerate}
\def\labelenumi{\arabic{enumi})}
\setcounter{enumi}{1}
\tightlist
\item
  设 \(\mathbf{T}\) 是一个完全正则空间，\(\mathfrak{F}\) 是 \(\mathcal{M}_{+}^{b}(\mathrm{T})\) 上的一个滤子，它紧收敛到一个有界测度 \(\mu\)。称 Borel 子集 \(\mathbf{A}\)（\(\mathbf{T}\) 的）为（对 \(\mathfrak{F}\) 的）收敛集，如果 \(\lim_{\lambda,\mathfrak{F}} \lambda_{\mathbf{A}} = \mu_{\mathbf{A}}\)。
\end{enumerate}

\begin{enumerate}
\def\labelenumi{\alph{enumi})}
\item
  若不相交的集合 \(\mathbf{A}_1\) 与 \(\mathbf{A}_2\) 都是收敛集，则 \(\mathbf{A}_1 \cup \mathbf{A}_2\) 也是收敛集。
\item
  设 A 是一个开集或闭集。为使 A 是收敛集，必要且充分的是 \(\lim_{\lambda,\mathfrak{F}} \lambda^{\bullet}(A) = \mu^{\bullet}(A)\)。若 A 是收敛集，则 \(T - A\) 也是收敛集。
\item
  设 A 是 T 中一个开或闭的收敛集，B 是一个收敛集且 T - B 也是收敛集。则 A ∪ B 与 A ∩ B 都是收敛集，它们的余集亦然。
\item
  由开收敛集生成的集环由收敛集组成。
\item
  设 \(\mathbf{A}\) 是 \(\mathbf{T}\) 的一个边界为 \(\mu\) -可忽略的子集。则 \(\mathbf{A}\) 是一个收敛集。
\item
  假设 T 是局部紧的或波兰的。证明 T 的每个紧子集 K 都包含在一个紧收敛集中（当 T 局部紧时，利用 Exer. 1 b)。若 T 是波兰空间，则利用同一习题构造 K 的一列有限覆盖 \(U_{p}\)（由 T 的开子集组成，边界 \(\mu\) -可忽略且直径 \(<2^{-p}\)）；证明 \(L=\bigcap_{p}\bigcup_{U\in\mathfrak{U}_{p}}\overline{U}\) 是紧的，并由 b) 断定它是一个收敛集）。
\end{enumerate}


\begin{enumerate}
\def\labelenumi{\alph{enumi})}
\setcounter{enumi}{6}
\tightlist
\item
  把 f) 的结果推广到完备度量空间上收敛测度序列的情形。
\end{enumerate}

\begin{enumerate}
\def\labelenumi{\arabic{enumi})}
\setcounter{enumi}{2}
\tightlist
\item
  设 \(\mathbf{T}\) 是一个波兰空间，\((\mu_n)\) 是 \(\mathbf{T}\) 上一列紧收敛到测度 \(\mu\) 的有界正测度。用 \(\mathfrak{C}\) 表示具有下列性质的子集 \(\mathbf{A}\)（\(\mathbf{T}\) 的）所成的集合：对每个 \(\varepsilon > 0\)，存在紧子集 \(\mathbf{K}\)（\(\mathbf{A}\) 的），使得 \(\sup_{n} \mu_n(\mathbf{A} - \mathbf{K}) < \varepsilon\)。用 \(\mathfrak{D}\) 表示子集 \(\mathbf{A}\)（\(\mathbf{T}\) 的，与其余集一起都属于 \(\mathfrak{C}\)）所成的集合。
\end{enumerate}

\begin{enumerate}
\def\labelenumi{\alph{enumi})}
\item
  若集合 A 与 A' 属于 C，则 \(A \cup A'\) 与 \(A \cap A'\) 也属于 C；由此推出 D 是 T 的子集所成的一个集环。
\item
  T 的每个边界为 \(\mu\) -可忽略的子集 A 都属于 \(\mathfrak{D}\)（对 A 的内部应用 Prokhorov 定理，并利用 Exer. 2 e)）。
\item
  每个 \(\mathbf{A} \in \mathfrak{D}\) 都是序列 \((\mu_n)\) 的收敛集（cf.~Exer. 2）；反之，每个开或闭的收敛集都属于 \(\mathfrak{D}\)。
\item
  设 \(\mathbf{A} \in \mathfrak{D}\)。证明存在一列两两不相交的紧集 \(\mathbf{K}_p\) 与子集 \(\mathbf{N}\)（\(\mathbf{T}\) 的），它们具有下列性质：\(\alpha\) ) 每个 \(\mathbf{K}_p\) 都是序列 \((\mu_n); \beta)\) 的收敛集；\(\mathbf{A} \subset \mathbf{N} \cup \bigcup_{p} \mathbf{K}_p\) 且 \(\bigcup_{p} \mathbf{K}_p \subset \mathbf{A} \cup \mathbf{N}; \gamma)\) 对每个 \(\varepsilon > 0\)，存在开邻域 \(\mathbf{U}\)（\(\mathbf{N}\) 的），使得 \(\sup_{n}\mu_n(\mathbf{U}) < \varepsilon\)（对 \(\mathbf{T}\) 的一个适当的子空间（它是一列开集的交）应用 Exer. 2 f) 与 Prokhorov 定理）。
\end{enumerate}

\begin{enumerate}
\def\labelenumi{\arabic{enumi})}
\setcounter{enumi}{3}
\item
  设 \(\mathbf{T}\) 是一个完全正则空间，\(t\) 是 \(\mathbf{T}\) 的一个点，\(\mathfrak{U}\) 是 \(\mathbf{T}\) 的一族 Borel 子集，它生成 \(t\) 的邻域滤子。设 \(\mathbf{I}\) 是带有一个滤子 \(\mathfrak{F}\) 的集合，\((\mu_i)_{i \in \mathbb{I}}\) 是 \(\mathbf{T}\) 上一族总质量为 1 的有界正测度。为使 \(\lim_{i, \mathfrak{F}} \mu_i = \varepsilon_t\)，必要且充分的是 \(\lim_{i, \mathfrak{F}} \mu_i(\mathrm{U}) = 1\)（对每个 \(\mathrm{U} \in \mathfrak{U}\)）。
\item
  设 E 是一个实 Hilbert 空间，具有一个标准正交基 \((x_{n})_{n\in \mathbf{N}}\)。设 T 是赋以弱化拓扑的空间 E，并设 \((a_{n})_{n\in \mathbf{N}}\) 是一列实数，使得 \(0 < a_{n} < 1\)（对所有 \(n\)），且 \(\lim_{n\to \infty}n^2\log a_n = 0\)。对每个 \(n\in \mathbf{N}\)，置 \(\mu_{n} = \sum_{p\in \mathbf{N}}a_{n}^{p}(1 - a_{n})\cdot \varepsilon_{n\cdot x_{p}}\)。证明 \(\mu_{n}\) 对每个 \(n \in \mathbf{N}\) 都是 T 上总质量为 1 的有界正测度，\(\lim_{n\to\infty}\mu_{n}=\varepsilon_{0}\)（紧收敛），并且 \(\lim_{n\to\infty}\mu_{n}(\mathrm{K})=0\)（对 T 的每个紧子集 \(\mathrm{K}\)）（对 0 的邻域组成的集合 \(\mathfrak{U}\)（形如 \(\{x \mid (a \mid x) \leqslant 1\}\)，其中 a 取遍 E）应用上一习题的判别法）。特别地，序列 \((\mu_{n})_{n \in \mathbf{N}}\) 的元素所成的集合是 \(\mathcal{M}_{+}^{b}(\mathbf{T})\) 中一个相对紧子集，但它不满足 Prokhorov 条件。
\end{enumerate}


\begin{enumerate}
\def\labelenumi{\arabic{enumi})}
\setcounter{enumi}{5}
\tightlist
\item
  设 \(\mathbf{T}\) 为一个拓扑空间，\(\mathbf{H}\) 为 \(\mathbf{T}\) 上正测度组成的一个集合；假设下列条件成立：
\end{enumerate}

\begin{enumerate}
\def\labelenumi{(\Alph{enumi})}
\setcounter{enumi}{21}
\tightlist
\item
  \(\mathbf{T}\) 的每个点都有一个开邻域 \(\mathbf{W}\)，使得 \(\sup_{\mu \in \mathbf{H}} \mu^{\bullet}(\mathbf{W})\) 有限。
\end{enumerate}

最后，设 \(\mathfrak{U}\) 为 H 上的一个超滤子。

\begin{enumerate}
\def\labelenumi{\alph{enumi})}
\item
  设 \(\mathbf{K}\) 是 \(\mathbf{T}\) 的一个紧子集。证明诱导测度 \(\mu_{\mathbf{K}} (\mu \in \mathbf{H})\) 关于 \(\mathfrak{U}\) 淡收敛到一个测度 \(\pi^{\mathbf{K}}\)（在 \(\mathbf{K}\) 上）。
\item
  设 \(\mathbf{K}\) 与 \(\mathbf{L}\) 是 \(\mathbf{T}\) 的两个紧子集，且 \(\mathbf{K} \subset \mathbf{L}\)；证明 \((\pi^{\mathbf{L}})_{\mathbf{K}} \geqslant \pi^{\mathbf{K}}\)。
\item
  对 T 的每个紧子集 K，用 \(i^{\mathbf{K}}\) 表示 K 到 T 中的典范单射。证明测度族 \(i^{\mathbf{K}}(\pi^{\mathbf{K}})\) 有一个上确界 \(\pi\)（在 \(\mathcal{M}_{+}(\mathrm{T})\) 中）。
\item
  设 \(f\) 是 \(\mathbf{T}\) 上的一个下半连续正函数，\(g\) 是 \(\mathbf{T}\) 上具有紧支撑的一个上半连续正函数。证明 \(\pi^{\bullet}(f) \leqslant \lim_{\mu, \mathfrak{U}} \mu^{\bullet}(f)\) 与 \(\pi^{\bullet}(g) \geqslant \lim_{\mu, \mathfrak{U}} \mu^{\bullet}(g)\)。
\end{enumerate}

¶ 7) 设 H 是拓扑空间 T 上有界正测度组成的一个集合。假设 \(\sup_{\mu\in H}\mu^{\bullet}(T)\) 有限，并且对每个 \(\varepsilon>0\)，存在 T 的一个紧子集 K

使得 \(\sup_{\mu \in \mathbf{H}}\mu^{\bullet}(\mathrm{T} - \mathrm{K}) < \varepsilon\)。证明上一习题的条件 (V) 是

满足的。设 \(\mathfrak{U}\) 为 H 上的一个超滤子。测度 \(\pi\) 按上一习题中的方式定义。

\begin{enumerate}
\def\labelenumi{\alph{enumi})}
\item
  证明 \(\pi^{\bullet}(g) \geqslant \lim_{\mu, \mathfrak{U}} \mu^{\bullet}(g)\) 对 T 上的每个上半连续有界正函数 \(g\) 成立。由此推出，\(\pi^{\bullet}(f) = \lim_{\mu, \mathfrak{U}} \mu^{\bullet}(f)\) 对 T 上每个有界函数 \(f\)（其不连续点集为 \(\pi\) -可忽略）成立。
\item
  当 T 完全正则时，由 a) 推出 H 对紧拓扑是相对紧的（这为 No.~5 的 Th. 1 在正测度的情形提供了一个新的证明）。
\end{enumerate}

¶ 8) 设 T 是一个完备可分度量空间，d 是它的距离。对 T 的每个闭子集 F 与每个实数 a > 0，用 \(F^{a}\) 表示 \(x \in T\) 中使 \(d(x, F) < a\) 的 x 所成的集合。给定 T 上两个有界正测度 \(\lambda\) 与 \(\mu\)，用 \(D(\lambda, \mu)\) 表示满足下列不等式的实数 a > 0 所成集合的下确界：对 T 的每个闭子集 F，\(\lambda^{\bullet}(F) \leqslant \mu^{\bullet}(F^{a}) + a\)，\(\mu^{\bullet}(F) \leqslant \lambda^{\bullet}(F^{a}) + a\)。

\begin{enumerate}
\def\labelenumi{\alph{enumi})}
\item
  证明 D 是 \(\mathcal{M}_+^b\) 上的一个距离，且 \(\mathrm{D}(\varepsilon_x,\varepsilon_y) = d(x,y)\) 对 T 的任意两个点 \(x\) 与 \(y\)（满足 \(d(x,y) < 1\)）成立。
\item
  设 \(f\) 是 \(T\) 到 \(T\) 中的一个 Borel 映射（cf.~TG, IX, §6, No.~3 或 GT, IX, §6, Exer. 16），\(\lambda\) 是 \(T\) 上的一个有界正测度。证明存在 T 上的一个有界正测度 \(\mu\)，使得对 T 的每个 Borel 子集 A 有 \(\mu^{\bullet}(\mathrm{A}) = \lambda^{\bullet}\big(f^{-1}(\mathrm{A})\big)\)。设 \(a > 0\)，并用 H 表示 \(x \in \mathrm{T}\) 中使 \(d(x, f(x)) \geqslant a\) 的 x 所成的集合；证明 H 是 T 中的 Borel 集，且 \(\mathrm{D}(\lambda, \mu) \leqslant \sup (a, \lambda(\mathrm{H}))\)（对 T 的每个闭子集 F，有 \(\mathrm{F} \cap (\mathrm{T} - \mathrm{H}) \subset f(\mathrm{F}^a)\) 且 \(\lambda^{\bullet}(\mathrm{F}) \leqslant \lambda^{\bullet}(\mathrm{F} \cap (\mathrm{T} - \mathrm{H})) + \lambda^{\bullet}(\mathrm{H})\)）。
\item
  设 \(\mu\) 与 \(\mu_n\)（对 \(n \geqslant 1\)）是 \(\mathbf{T}\) 上的有界正测度，使得 \(\lim_{n \to \infty} \mathrm{D}(\mu_n, \mu) = 0\)。证明序列 \((\mu_n)\) 紧收敛到 \(\mu\)（证明 \(\lim_{n \to \infty} \mu_n^\bullet(\mathrm{T}) = \mu^\bullet(\mathrm{T})\)，且 \(\mu^\bullet(\mathrm{F}) \geqslant \limsup_{n \to \infty} \mu_n^\bullet(\mathrm{F})\)（对每个闭子集 \(\mathbf{F}\)（\(\mathbf{T}\) 的））；再按 §2, No.~6 的 Lemma 3 的方法推出 \(\mu^\bullet(f) \leqslant \liminf_{n \to \infty} \mu_n^\bullet(f)\)（对每个下半连续函数 \(f \geqslant 0\)））。
\item
  反之，证明对每个有界正测度序列 \(\mu_n\)（紧收敛到 \(\mu\)，在 \(\mathcal{M}_+^b\) 中），有 \(\lim_{n \to \infty} D(\mu_n, \mu) = 0\)。可以按如下方式进行：设 \(\varepsilon > 0\)，并设 \(K\) 是一个紧子集（\(T\) 的），使得 \(\mu^\bullet(T - K) < \varepsilon\) 且 \(\sup_{n} \mu_n^\bullet(T - K) < \varepsilon\)；
\end{enumerate}

构造一个有限族 \((\mathrm{B}_{i})_{1\leqslant i\leqslant p}\)（由边界 \(\mu\) -可忽略、直径 \(\leqslant\varepsilon/2\)、两两不相交且并为包含 K 的集合组成）（cf.~Exer. 1）。设 f 是 T 到 T 中的一个映射，它在每个集合 \(B_{1},\ldots,B_{p}\) 与 \(\mathbb{C}(B_{1}\cup\cdots\cup B_{p})\) 上取常值，且 \(f(\mathrm{B}_i) \subset \mathrm{B}_i\)（对 \(1 \leqslant i \leqslant p\)）。用下列条件定义测度 \(\pi_n\) 与 \(\pi\)：对 T 的每个 Borel 子集 A，\(\pi_n^\bullet(\mathrm{A}) = \mu_n^\bullet(\overline{f}(\mathrm{A}))\) 且 \(\pi^\bullet(\mathrm{A}) = \mu^\bullet(f^{-1}(\mathrm{A}))\)；由 b) 推出关系式 \(\mathrm{D}(\pi_n, \mu_n) \leqslant \varepsilon\)，\(\mathrm{D}(\pi, \mu) \leqslant \varepsilon\)，并证明 \(\lim_{n \to \infty} \mathrm{D}(\pi_n, \pi) = 0\)。

\begin{enumerate}
\def\labelenumi{\alph{enumi})}
\setcounter{enumi}{4}
\tightlist
\item
  \(M_{+}^{b}\) 上的距离 D 与紧拓扑相容（注意 \(M_{+}^{b}\) 对紧拓扑是可度量化的）。
\end{enumerate}

¶ 9) 记号与假设同上一习题。设 \((\mu_n)\) 是 \(\mathcal{M}_+^b\) 中对距离 D 的一个 Cauchy 序列；证明序列 \((\mu_n)\) 是收敛的（这给出了空间 \(\mathcal{M}_+^b\) 对紧拓扑是波兰空间这一事实的一个新证明）。可以按如下方式进行：

\begin{enumerate}
\def\labelenumi{\alph{enumi})}
\item
  设 \(\varepsilon > 0\) 与 \(a > 0\) 是两个实数；存在一个有限子集 \(F\)（\(T\) 的），使得 \(\sup_{n} \mu_n^\bullet(T - F^a) \leqslant \varepsilon\)（选取一个整数 \(N \geqslant 1\) 与一个紧子集 \(K\)（\(T\) 的），使得 \(\sup_{n \geqslant N} D(\mu_n, \mu_N) \leqslant \varepsilon/2\) 且 \(\mu_N^\bullet(T - K) \leqslant \varepsilon/2\)；由此推出 \(\sup_{n \geqslant N} |\mu_n^\bullet(T) - \mu^\bullet(T)| \leqslant \varepsilon\) 与 \(\sup_{n \geqslant N} \mu_n^\bullet(T - K^{a/2}) \leqslant \varepsilon\)；最后选取一个有限子集 \(F\)（\(T\) 的），使得 \(K^{a/2} \subset F^a\) 且 \(\sup_{n \geqslant N} \mu_n^\bullet(T - F^a) \leqslant \varepsilon\)）。
\item
  设 \(\varepsilon > 0\)；存在一个紧子集 \(K\)（\(T\) 的），使得 \(\sup_{n} \mu_n^\bullet(T - K) \leqslant \varepsilon\)（对每个整数 \(p \geqslant 1\)，选取一个有限子集 \(F_p\)（\(T\) 的），使得 \(\mu_n^\bullet(T - (F_p)^{2^{-p}}) \leqslant \varepsilon / 2^p\)，并置 \(K = \bigcap_{p\geqslant 1}(\overline{(F_p)^{2^{-p}}})\)。
\item
  由 \(b\) ) 推出，Cauchy 序列 \((\mu_n)\) 有一个收敛的子序列（对距离 D）。
\end{enumerate}

¶ 10) 设 T 是一个完全正则空间；用 S 表示 T 上满足下列性质的实函数 \(f \geqslant 1\) 所成的集合：对每个实数 c，使 \(f(t) \leqslant c\) 的 T 的点 t 所成的集合是紧的。~对每个 \(f \in S\)，用 \(\mathcal{M}_f\) 表示 T 上有界测度 \(\mu\)（满足 \(|\mu|^{\bullet}(f) \leqslant 1\)）所成的集合。

\begin{enumerate}
\def\labelenumi{\alph{enumi})}
\item
  为使 \(\mathcal{M}^{b}(T)\) 的子集 H 满足 Prokhorov 条件，必要且充分的是存在 \(f \in S\)，使得 \(H \subset M_{f}\)（对必要性，取 f 形如 \(c(1 + \sum_{n \geqslant 1} n \cdot f_{n})\)，其中 \(f_{n}\) 是这样的集合 \(U_{n}\) 的特征函数：\(\mathrm{T} - \mathrm{U}_n\) 是紧的，且 \(\sup_{\mu \in \mathbf{H}} |\mu|^{\bullet} (\mathrm{U}_n) \leqslant 2^{-n}\)。
\item
  假设 T 是局部紧的。为使 \(\mathcal{M}^{b}(T)\) 的子集 H 是相对紧的（对紧拓扑），必要且充分的是存在连续函数 \(f \in S\)，使得 \(H \subset M_{f}\)。
\item
  设 C 是 \(\mathcal{M}_{+}^{b}(\mathrm{T})\) 中一个闭凸锥。证明 \(\mathrm{C} \cap \mathcal{M}_{f}\) 对每个 \(f \in \mathcal{S}\) 都是 C 的一个帽（TVS, II, §7, No.~2, Def. 3），并由此推出 C 是它的帽的并。用 E 表示 C 的极端生成元的并。假设 T 是 Souslin 空间。证明对每个 \(\pi \in \mathrm{C}\)，存在 \(\mathcal{M}_{+}^{b}(\mathrm{T})\) 中包含于 E 的一个 Borel 子集 B 与 B 上一个总质量为 1 的正测度 P，使得 \(\pi = \int_{\mathrm{B}} \mu d\mathrm{P}(\mu)\)（应用 Choquet 积分表示定理）。
\end{enumerate}

\begin{enumerate}
\def\labelenumi{\arabic{enumi})}
\setcounter{enumi}{10}
\item
  设 T 是一个完全正则空间。假设对每个紧空间 K 与每个 T 到 K 中的连续映射 f，都给定 K 上的一个正测度 \(\mu_{f,K}\)；假设 \(g(\mu_{f,K}) = \mu_{g \circ f,L}\)（对任意紧空间 K 与 L 以及任意连续映射 \(f : T \to K\) 与 \(g : K \to L\)）；再假设对任意 f 与 K，测度 \(\mu_{f,K}\) 集中在 \(f(T)\) 上。证明存在唯一的 T 上的有界正测度 \(\pi\)，使得对任意 f 与 K 有 \(\mu_{f,K} = f(\pi)\)（利用 T 的 Stone--Čech 紧化的泛性质）。
\item
  设 T 是一个完全正则空间，是一个局部紧空间 L 的子空间。对 T 上的每个测度 \(\mu\)（无论正与否），存在 L 的一个包含 T 的局部紧子空间 \(T'\)，以及一个集中在 T 上的测度 \(\mu'\)（在 \(T'\) 上），它在 T 上诱导出 \(\mu\)。
\end{enumerate}

¶ 13) *设 G 是一个局部紧交换群。用 \(\widehat{G}\) 表示 G 的对偶群，用 \(\alpha\) 表示 \(\widehat{G}\) 上的一个 Haar 测度。对 G 上的每个有界测度 \(\mu\)，用 \(\mathcal{F}\mu\) 表示 \(\mu\) 的 Fourier 变换，它是群 \(\widehat{G}\) 上的一个函数（cf.~Theor. spec., Ch. II, §1, No.~2）。

\begin{enumerate}
\def\labelenumi{\alph{enumi})}
\item
  Fourier 变换 \(\mathcal{F}\) 是从赋以紧拓扑的 \(\mathcal{M}_{+}^{b}(\mathrm{G})\) 到赋以紧收敛拓扑的 \(\mathcal{C}^b (\widehat{\mathrm{G}})\) 中的一个连续映射。（利用 No.~6 的 Prop. 13 的推论。）
\item
  设 \(\mathfrak{F}\) 是 \(\mathcal{M}_+^b (\mathrm{G})\) 上的一个滤子，\(\Phi\) 是 \(\widehat{\mathrm{G}}\) 上的一个有界连续函数；假设 \(\lim_{\lambda ,\mathfrak{F}}(\mathcal{F}\lambda)\cdot \alpha = \Phi \cdot \alpha\)（淡收敛）且 \(\lim_{\lambda ,\mathfrak{F}}(\mathcal{F}\lambda)(0) = \Phi (0)\)。证明滤子 \(\mathfrak{F}\) 紧收敛到一个测度 \(\mu\)，使得 \(\mathcal{F}\mu = \Phi\)（应用 Stone-Weierstrass 定理证明：G 上具有紧支撑的连续函数的 Fourier 变换所成的集合，是 \(\mathcal{C}^0 (\overline{\mathrm{G}})\) 中对一致收敛的稠密线性子空间；接着注意 \(\lim_{\lambda ,\mathfrak{F}}\int (\mathcal{F}\lambda)\cdot u d\alpha = \int \Phi u d\alpha\) 对每个 \(u\in \mathrm{E}\) 成立，且滤子 \(\mathfrak{F}\) 含有有界测度的一个淡紧集合；由此推出滤子 \(\mathfrak{F}\) 至多有一个聚点，然后推出它紧收敛）。
\item
  设 \(\mathfrak{F}\) 是 \(\mathcal{M}_{+}^{b}(\mathrm{G})\) 上具有可数基的一个滤子。为使 \(\mathfrak{F}\) 紧收敛到一个有界正测度 \(\mu\)，必要且充分的是存在一个连续函数 \(\Phi\)（在 \(\widehat{\mathbf{G}}\) 上），使得 \(\mathcal{F}\lambda\) 逐点收敛到 \(\Phi\)（关于 \(\mathfrak{F}\)），此时 \(\mathcal{F}\mu = \Phi\)（\(P\) . Lévy 定理）。*
\end{enumerate}

\exercisesection{}

¶ 1) 对 R 的每个紧区间 K，用 \(\mathcal{C}(K)\) 表示 K 上连续函数的实向量空间，赋以一致收敛拓扑；用 \(\mathcal{D}(K)\) 表示 R 上无穷可微且在 K 之外为零的实函数所成的向量空间。给 \(\mathcal{D}(K)\) 赋以使映射 \(f \mapsto D^{p}f|K\)（对每个正整数 \(p\)）都连续的最粗拓扑（D 是求导算子）。置 \(\mathcal{D} = \bigcup_{K} \mathcal{D}(K)\)，并给这个空间赋以这样的局部凸拓扑：它是子空间

\(\mathcal{D}(\mathbf{K})\) 的诸拓扑的直极限。

\begin{enumerate}
\def\labelenumi{\alph{enumi})}
\tightlist
\item
  对每个整数 \(p \geqslant 0\) 与每个函数 \(f \in \mathcal{D}(\mathbf{K})\)，置
\end{enumerate}

\[
\mathrm{Q} _ {p} (f) = \int_ {\mathrm{K}} \left(\mathrm{D} ^ {p} f (x)\right) ^ {2} d x.
\]

证明 \(\mathbf{Q}_p\) 是 \(\mathcal{D}(\mathbf{K})\) 上的一个正二次型，范数序列 \(\mathbf{Q}_p^{1/2}\) 定义了 \(\mathcal{D}(\mathbf{K})\) 的拓扑，且 \(\operatorname{Tr}(\mathbf{Q}_{p+1}/\mathbf{Q}_p)\)（对每个 \(p \geqslant 0\)）有限。（对每个实数 \(t\)，用 \(\mathrm{I}_t\) 表示区间 \([t, +\infty[\) 的特征函数；利用公式 \(\mathrm{D}^p f(t) = \int_{\mathbf{K}} \mathrm{D}^{p+1} f \cdot \mathrm{I}_t dx\) 与 Bessel 不等式，证明 \(\sum_{i=1}^{n} \mathbf{Q}_p(f_i) \leqslant l^2\) 对每个有限序列 \(f_1, \ldots, f_n\)（由属于 \(\mathcal{D}(\mathbf{K})\) 的函数组成，对 \(\mathbf{Q}_{p+1}\) 标准正交）成立；数 \(l\) 是区间 K 的长度。）由此推出 \(\mathcal{D}(\mathbf{K})\) 是一个核空间。

\begin{enumerate}
\def\labelenumi{\alph{enumi})}
\setcounter{enumi}{1}
\tightlist
\item
  证明空间 \(\mathcal{D}\) 是核的。（先确立 \(\mathcal{D}\) 中一个非零函数的存在性，并由此推出函数 \(h \geqslant 0\)（\(\mathcal{D}\) 中的）的存在性，使得 \(\sum_{n \in \mathbf{Z}} h(x - n) = 1\)（对每个 \(x \in \mathbf{R}\)）。设 \(K\) 是 \(\mathbf{R}\) 的一个紧区间，使得 \(h\) 在
\end{enumerate}

K 之外为零；对每个整数 \(n\)，置 \(h_n(x) = h(x - n)\) 与 \(\mathrm{K}_n = \mathrm{K} + n\)。设 V 是 \(\mathcal{D}\) 中 0 的一个凸邻域；存在正整数 \(p_n\)，使得每个函数 \(f \in \mathcal{D}(\mathrm{K}_n)\)（满足 \(\mathrm{Q}_{p_n}(f) \leqslant 1\)）都属于 V。用下列公式定义 \(\mathcal{D}\) 上的连续二次型 Q 与 R：\(\mathrm{Q}(f) = \sum_{n \in \mathbf{Z}} 2^{2n} \mathrm{Q}_{p_n}(f \cdot h_n)\) 与 \(\mathrm{R}(f) = \sum_{n \in \mathbf{Z}} 2^{3n} \mathrm{Q}_{p_n + 1}(f \cdot h_n)\)；证明 V 包含由 \(f \in \mathcal{D}\)（满足 \(\mathrm{Q}(f) \leqslant 1\)）组成的集合，且 \(\operatorname{Tr}(\mathrm{R}/\mathrm{Q})\) 有限。）

\(^{*}\) c) 把上述结果推广到 \(R^{n}\) 上具有紧支撑的无穷可微函数。\(^{*}\)

\specsection{附录的习题}

\begin{enumerate}
\def\labelenumi{\arabic{enumi})}
\tightlist
\item
  设 E 是特征 \(\neq 2\) 的交换域上的一个有限维向量空间。用 H 表示 E 上一个非退化二次型，用 S 表示 \(E \times E\) 上使得对 E 中所有 x 有 \(\mathrm{H}(x) = \mathrm{S}(x, x)\) 的对称双线性型。设 Q 是 E 上的一个二次型。
\end{enumerate}

\begin{enumerate}
\def\labelenumi{\alph{enumi})}
\item
  存在自同态 \(u\)（\(\mathbf{E}\) 的），它由关系式 \(\mathbf{Q}(x) = \mathrm{S}(u(x), x)\) 与 \(\mathrm{S}(u(x), y) = \mathrm{S}(x, u(y))\)（\(x\)、\(y\) 在 \(\mathbf{E}\) 中）刻画。置 \(\operatorname{Tr}(\mathbf{Q}/\mathbf{H}) = \operatorname{Tr}(u)\)。
\item
  推广 No.~1 的 Remark 3。
\item
  若 \((e_i)_{1\leqslant i\leqslant m}\) 是 E 的对 H 标准正交的基，则 \(\operatorname{Tr}(\mathbf{Q}/\mathbf{H}) = \sum_{i=1}^{m}\mathrm{Q}(e_i)\)。
\end{enumerate}

\begin{enumerate}
\def\labelenumi{\arabic{enumi})}
\setcounter{enumi}{1}
\tightlist
\item
  设 E 是一个实 Hilbert 空间，Q 是 E 上一个连续正二次型，H 是 E 上的二次型 \(x \mapsto \|x\|^2\)。
\end{enumerate}

\begin{enumerate}
\def\labelenumi{\alph{enumi})}
\item
  证明 \(\operatorname{Tr}(\mathbf{Q}/\mathbf{H}) = \sum_{i \in I} \mathbf{Q}(e_i)\) 对 E 的每个标准正交基 \((e_i)_{i \in I}\) 成立。（设 \((a_1, \ldots, a_p)\) 是 E 中一个有限标准正交族；对每个 \(\varepsilon > 0\)，存在 I 的有限子集 J 与元素 \(a_1', \ldots, a_p'\)，它们是 \(e_i\)（对 \(i \in J\)）的线性组合，且满足 \(\|a_j - a_j' \| < \varepsilon\)（对 \(1 \leqslant j \leqslant p\)）；于是 \(\sum_{j=1}^{p} \mathbf{Q}(a_j') \leqslant \sum_{i \in J} \mathbf{Q}(e_j) \leqslant \sum_{i \in I} \mathbf{Q}(e_i)\)。由此推出 \(\sum_{j=1}^{p} \mathbf{Q}(a_j) \leqslant \sum_{i \in I} \mathbf{Q}(e_i)\)。）
\item
  由 \(a\) ) 推出 Prop. 3 的一个新证明。
\item
  设 \(\mathbf{E}_0\) 是 \(\mathbf{E}\) 的一个稠密线性子空间，\(\mathbf{Q}_0\)（相应地，\(\mathbf{H}_0\)）是 \(\mathbf{Q}\)（相应地，\(\mathbf{H}\)）在 \(\mathbf{E}_0\) 上的限制。证明 \(\operatorname{Tr}(\mathbf{Q}/\mathbf{H}) = \operatorname{Tr}(\mathbf{Q}_0/\mathbf{H}_0)\)。（设 \((e_1, \ldots, e_n)\) 是 \(\mathbf{E}\) 中一个线性无关的序列，生成一个子空间 \(\mathbf{F}\)。用 \(\mathbf{Q}_{\mathbf{F}}\)（相应地，\(\mathbf{H}_{\mathbf{F}}\)）表示 \(\mathbf{Q}\)（相应地，\(\mathbf{H}\)）在 \(\mathbf{F}\) 上的限制。利用 No.~1 的 Remark 3，证明 \(\operatorname{Tr}(\mathbf{Q}_{\mathbf{F}}/\mathbf{H}_{\mathbf{F}})\) 是 \((e_1, \ldots, e_n)\) 的连续函数。）
\end{enumerate}

\begin{enumerate}
\def\labelenumi{\arabic{enumi})}
\setcounter{enumi}{2}
\item
  设 E 与 F 是两个实向量空间，u 是 E 到 F 中的一个线性映射；若 Q 与 H 是 F 上的正二次型，使得 \(\mathrm{H}(x)=0\) 蕴含 \(\mathrm{Q}(x)=0\)（对 \(x\in F\)），则 \(\mathrm{Tr}\left(\mathrm{Q}\circ\mathrm{u}/\mathrm{H}\circ\mathrm{u}\right)\leqslant\mathrm{Tr}\left(\mathrm{Q}/\mathrm{H}\right)\)。
\item
  设 I 是一个集合，\(l^2(\mathbf{I})\) 是由实数族 \(\mathbf{x} = (x_i)_{i \in \mathbf{I}}\)（满足 \(\sum_{i \in \mathbf{I}} x_i^2\) 有限）组成的向量空间。设 \((\lambda_i)_{i \in \mathbf{I}}\) 是一个可和的正数族。对每个 \(\mathbf{x}\)（在 \(l^2(\mathbf{I})\) 中），置 \(Q(\mathbf{x}) = \sum_{i \in \mathbf{I}} \lambda_i x_i^2\) 与 \(H(\mathbf{x}) = \sum_{i \in \mathbf{I}} x_i^2\)。证明 \(\operatorname{Tr}(Q / H) = \sum_{i \in \mathbf{I}} \lambda_i\)。（利用 Exer. 2 a)。
\item
  设 E 是一个实向量空间，\((\lambda_i)_{i \in I}\) 是一个可和的正数族，\((y_i)_{i \in I}\) 是 E 上线性形式组成的一个族。对每个 \(x \in E\)，置 \(\mathrm{H}(x) = \sum_{i \in I} \langle x, y_i \rangle^2\) 与 \(\mathrm{Q}(x) = \sum_{i \in I} \lambda_i \langle x, y_i \rangle^2\)；假设 \(\mathrm{H}(x)\)（对每个 \(x \in E\)）有限。证明 \(\mathrm{Q}(x)\)（对每个 \(x \in E\)）有限，Q 与 H 是 E 上的正二次型，且 \(\operatorname{Tr}(\mathrm{Q}/\mathrm{H}) \leqslant \sum_{i \in I} \lambda_i\)。（置 \(u(x) = (\langle x, y_i \rangle)_{i \in I}\)，并对线性映射 \(u\)（E 到 \(l_2(I)\) 中的）应用 Exer. 3。）
\item
  设 E 是一个实向量空间，Q、H 与 \(H_{0}\) 是 E 上的正二次型。假设 \(H \leqslant a \cdot H_{0}\)，其中 a 是一个正实数。证明 \(\operatorname{Tr}(Q/H_{0}) \leqslant a \cdot \operatorname{Tr}(Q/H)\)。（利用 No.~1 的 Remark 1 化为 E 有限维的情形，然后利用 E 的一个对 H 与 \(H_{0}\) 都正交的基的存在性，由 Prop. 1 得出结论。）
\item
  设 E 是一个实 Hilbert 空间。证明 E 上的 Sazonov 拓扑由半范数 \(Q^{1/2}\) 组成的集合定义，其中 Q 取遍 E 上核正二次型所成的集合。（利用 Exer. 6。）
\item
  设 E 与 F 是两个实 Hilbert 空间。在 \(E \otimes F\) 上存在一个 Hausdorff 准 Hilbert 空间结构，使得 \((x_{1} \otimes y_{1}|x_{2} \otimes y_{2}) = (x_{1}|x_{2}) \cdot (y_{1}|y_{2})\)（对 E 中的 \(x_{1}, x_{2}\) 与 F 中的 \(y_{1}, y_{2}\)）。用 \(E \otimes_{2} F\) 表示 \(E \otimes F\) 的 Hilbert 空间完备化。
\end{enumerate}

\begin{enumerate}
\def\labelenumi{\alph{enumi})}
\item
  若 \(\mathbf{E}'\)（相应地，\(\mathbf{F}'\)）是 \(\mathbf{E}\)（相应地，\(\mathbf{F}\)）的一个闭线性子空间，证明 \(\mathbf{E}' \otimes_2 \mathbf{F}'\) 可以等同于 \(\mathbf{E} \otimes_2 \mathbf{F}\) 中由元素 \(x \otimes y\)（满足 \(x \in \mathbf{E}'\) 与 \(y \in \mathbf{F}'\)）生成的闭线性子空间。
\item
  假设 E（相应地，F）是闭线性子空间族 \((\mathrm{E}_{\alpha})_{\alpha\in\mathbf{A}}\)（相应地，\((\mathrm{F}_{\beta})_{\beta\in\mathbf{B}}\)）的 Hilbert 和（TVS, V, §2, No.~2, Def. 2）。证明 \(E\otimes_{2}F\) 是闭线性子空间族 \((\mathrm{E}_{\alpha}\otimes_{2}\mathrm{F}_{\beta})_{(\alpha,\beta)\in\mathbf{A}\times\mathbf{B}}\) 的 Hilbert 和。
\item
  若 \((e_i)_{i\in \mathbf{I}}\)（相应地，\((f_j)_{j\in \mathbf{J}})\)）是 E（相应地，F）的标准正交基，则族 \((e_i\otimes f_j)_{(i,j)\in \mathbf{I}\times \mathbf{J}}\) 是 \(\mathbf{E}\otimes_2\mathbf{F}\) 的标准正交基。
\item
  设 \(\mathbf{G}\) 是一个实 Hilbert 空间。定义从 \(\mathbf{E} \otimes_{2} \mathbf{F}\) 到 \(\mathbf{F} \otimes_{2} \mathbf{E}\) 上以及从 \((\mathbf{E} \otimes_{2} \mathbf{F}) \otimes_{2} \mathbf{G}\) 到 \(\mathbf{E} \otimes_{2} (\mathbf{F} \otimes_{2} \mathbf{G})\) 上的典范同构。
\item
  设 \(E_{1}\) 与 \(F_{1}\) 是两个实 Hilbert 空间，\(u:E\to E_{1}, v:F\to F_{1}\) 是连续线性映射。证明 \(u\otimes v\) 可以由连续性延拓为连续线性映射 \(u\otimes_{2}v\)（\(E\otimes_{2}F\) 到 \(E_{1}\otimes_{2}F_{1}\) 中的），且 \(\|u\otimes_{2}v\|=\|u\|\cdot\|v\|\)。
\end{enumerate}

\begin{enumerate}
\def\labelenumi{\arabic{enumi})}
\setcounter{enumi}{8}
\tightlist
\item
  设 E 与 F 是两个实 Hilbert 空间。
\end{enumerate}

\begin{enumerate}
\def\labelenumi{\alph{enumi})}
\item
  证明存在一个连续线性映射 \(\varphi\)（\(\mathbf{E} \otimes_2 \mathbf{F}\) 到 \(\mathcal{L}(\mathbf{E}; \mathbf{F})\) 中的），它由下列条件刻画：\(\varphi(x \otimes y)(x') = (x|x') \cdot y\)（对 \(x, x'\)（\(\mathbf{E}\) 中的）与 \(y\)（\(\mathbf{F}\) 中的））。证明 \(\varphi\) 的范数为 1，且 \(\varphi(\mathbf{E} \otimes \mathbf{F})\) 是 \(\mathbf{E}\) 到 \(\mathbf{F}\) 中的有限秩连续线性映射所成的集合。
\item
  证明 \(\varphi\) 是 \(E \otimes_2 F\) 到 E 到 F 中的 Hilbert--Schmidt 线性映射所成的集合 HS (E, F) 上的一个双射（利用 Exer. 8 c) 与 Prop. 3）。给 HS (E, F) 赋以从 \(E \otimes_2 F\) 的 Hilbert 空间结构借助双射 \(\varphi\) 导出的结构；相应的范数记作 \(\|u\|_2\)。设 \(u \in \mathrm{HS}(E, F)\)；定义 E 上的正二次型 \(Q_u\) 与 H 如下：\(Q_u(x) = \|u(x)\|^2\) 与 \(H(x) = \|x\|^2\)。证明 \(\|u\|_2^2 = \operatorname{Tr}(Q_u / H)\)。
\item
  设 \(\mathbf{E}_1\) 与 \(\mathbf{F}_1\) 是两个实 Hilbert 空间，\(u: \mathbf{E}_1 \to \mathbf{E}\)、\(v: \mathbf{F} \to \mathbf{F}_1\) 是连续线性映射。设 \(\varphi_1\) 是 \(\mathbf{E}_1 \otimes_2 \mathbf{F}_1\) 到 \(\mathrm{HS}(\mathbf{E}_1, \mathbf{F}_1)\) 上的按类似于 \(\varphi\) 的方式定义的同构。证明 \(v \circ \varphi(t) \circ u = \varphi_1((u^* \otimes v)(t))\) 对所有 \(t \in \mathbf{E} \otimes_2 \mathbf{F}\) 成立。由此推出，对每个 Hilbert-Schmidt 映射 \(w\)（\(\mathbf{E}\) 到 \(\mathbf{F}\) 中的），线性映射 \(v \circ w \circ u\)（\(\mathbf{E}_1\) 到 \(\mathbf{F}_1\) 中的）是 Hilbert-Schmidt 的，且 \(\|v \circ w \circ u\|_2 \leqslant \|v\| \cdot \|w\|_2 \cdot \|u\|\)。
\end{enumerate}

\specialchapter{历史注记}

（注意 --- 罗马数字指本注记末尾的参考文献。）

尽管拓扑学与测度论之间联系的研究可以追溯到现代实变量函数论的开端，但在 Hausdorff 拓扑空间中的积分以一般的方式得到彻底研究，还只是最近的事。在叙述先于本综合工作的那些工作的历史之前，我们回顾关于拓扑与测度之间关系的观念演化的几个阶段。

对 Lebesgue 而言，问题只在于积分一个或几个实变量的函数。1913 年，Radon 定义了 \(\mathbf{R}^n\) 上的一般测度以及相应的积分；这一理论在 Ch. de la Vallée Poussin 的书 (I) 中有详细的阐述，并且自始至终依赖于 Euclid 空间的拓扑性质。稍后，在 1915 年，Fréchet 在 (II, \(a\) ) 中定义了赋有一个 σ-集环的集合上的“抽象”测度以及关于这些测度的积分；他注意到，用这种方式可以不借助拓扑方法而建立 Lebesgue 理论的主要结果。他用取自 (II, \(b\) ) 引言（第一部分）的下面这段话来为自己的工作辩护：« Que par exemple dans l'espace à une infinité de coordonnées où diverses applications de l'Analyse avaient conduit à diverses définitions non équivalentes d'une suite convergente, on remplace une de ces définitions par une autre, rien ne sera changé dans les propriétés des familles et fonctions additives d'ensembles dans ces espaces».(*) Fréchet 的研究由 Carathéodory 完成，集函数扩张为测度的一个重要定理应归功于后者。Saks 的书 (III) 的开头对这一观点作了扼要的阐述。

Haar 测度在局部紧群上的发现（cf.~第七章与第八章的历史注记）以及它随即获得的众多应用，随后 Weil 与 Gelfand 在调和分析方面的工作，大约在 1940 年前后导致了这一观点的深刻改变：在这类问题中，把测度视为连续函数空间上的线性形式最为方便。这一方法使人不得不限于紧空间或局部紧空间，但对几乎所有的应用而言这并无不便：更好的是，J. Tate 与 A. Weil 在 p-adic 群与 adèle 群上引入调和分析，使解析数论获得了惊人的更新。

正是从一个完全不同的方向，产生了通过考虑非局部紧拓扑空间上的测度来扩大这一观点的需要：概率论逐渐引导到对这类空间的研究，并提供了许多非平凡的例子。这些发展对测度论的影响之所以姗姗来迟，其原因或许在于概率论的相对孤立，直到不久以前它还处于传统数学学科的边缘。

\specsection{序列空间上的测度}

经典概率论发展最充分的分支之一是极限定理的分支（大数定律、趋于 Gauss--Laplace 律，\ldots)；这涉及对由牵涉极大总体的现象所显示的统计规律性概念的深化。这些问题的正确数学表述需要在序列空间上引入测度；这些空间构成有限维空间最明显的推广，是 Fréchet、Lévy、Lusin\ldots{} 在 1920 年前后所从事的“广义分析”研究的首选对象。新概率论方法的创立者 Khintchine 与 Kolmogoroff 都是 Lusin 的学生，而 Lévy 很快就转向了概率问题，这也并非偶然：这些构成了新方法的试金石。

序列空间上的测度第一次隐含地出现于 E. Borel 1909 年讨论可数概率的著作 (IV) 中。Borel 的一个高度独创的想法，是把他刚得到的概率结果应用于证明几乎每个介于 0 与 1 之间的实数的十进制展开所具有的性质。这一应用基于下列基本注记：用给定基 \(q\)（\(q \geqslant 2\)）下展开式中的数字（或“数码”）的序列来定义每个介于 0 与 1 之间的实数；如果依次随机地抽取一个数 \(x\) 的各个数字，以相互独立、相同的概率 \(1 / q\)（对 \(0,1,\ldots ,q - 1\)），则 \(x\) 落在 {[}0,1{[} 的一个区间中的概率等于该区间的长度。

1923 年，Steinhaus (V) 严格地建立了这些结果，并描述了 Borel 所考虑的无穷随机抽取序列的精确数学模型：为简单起见取 q=2，并用 I 表示两个元素的集合 \(\{0,1\}\)；给 I 装备测度 \(\mu\)，它由 \(\mu(0)=\mu(1)=\frac{1}{2}\) 定义；乘积空间 \(I^{N}\) 的元素是由等于 0 或 1 的数组成的序列 \(\varepsilon=(\varepsilon(n))_{n\in\mathbf{N}}\)，而映射 \(\varphi:\varepsilon\mapsto\sum_{n\geqslant0}\varepsilon(n)\cdot2^{-n-1}\) 除一个可数集外是 \(I^{N}\) 到区间 {[}0,1{]} 上的一个双射；此外，\(\varphi^{-1}\) 把 {[}0,1{]} 上的 Lebesgue 测度变换为 \(I^{N}\) 上的测度 P，即各因子上测度 \(\mu\) 的乘积。实际上，Steinhaus 当时并没有乘积测度的构造可用；他利用拟双射 \(\varphi\) 的存在性，构造 \(I^{N}\) 上的测度 P（从 {[}0,1{]} 上的 Lebesgue 测度出发），然后给出 P 的一个公理化刻画。这样得到的同构使人们能够把概率的语言翻译为测度的语言，并应用关于 Lebesgue 积分的已知定理。

在同一著作中，Steinhaus 考虑了随机级数 \(\sum_{n\geqslant 0}\sigma_n\cdot a_n\)，其中符号 \(\sigma_{n} = \pm 1\) 相互独立且以相同的概率 \(\frac{1}{2}\) 随机选取；在 1928 年至 1935 年之间，他研究了许多其他的随机级数。另一方面，Paley、Wiener 与 Zygmund 考虑了形如 \(\sum_{n = -\infty}^{\infty}a_n\exp (2\pi i(nt + \Phi_n))\) 的随机 Fourier 级数(1)；其中“振幅”\(a_{n}\) 是固定的，而“相位”\(\Phi_{n}\) 是在 {[}0,1{]} 上均匀分布的独立随机变量。虽然从一个问题到另一个问题解析上的困难相差悬殊，但用测度论语言所作的翻译在所有情形中都是相同的，并且是 Borel 与 Steinhaus 所处理情形的一个推广；问题在于构造 \(\mathbf{R}^{\mathbf{N}}\) 上的一个测度，它是一个测度族的乘积，族中的测度都等同于同一个质量为 1 的正测度 \(\mu\)（\(\mathbf{R}\) 上的）；例如，上述随机 Fourier 级数对应于 \(\mu\) 是 {[}0,1{]} 上的 Lebesgue 测度的情形。

为构造这样的乘积测度，可以用两种方法。第一种是直接方法，由 Daniell (VI, a)) 于 1918 年首次精确阐述；它在 1934 年被 Jessen (VII) 重新发现，后者对 \(\mu\) 是 {[}0, 1{]} 上 Lebesgue 测度的情形作了详细研究。第二种方法是寻求类似于 Steinhaus 的技巧，以化归到 {[}0, 1{]} 上的 Lebesgue 测度；在还没有一般测度论的完整阐述可供使用时，这种做法有方便的优点，因为它允许使用 Lebesgue 的定理而无需新的证明。\(^{(2)}\)

\specsection{Brown 运动的理论}

这一理论在当代科学发展中占有特殊的地位，它见证了物理问题与“纯”数学之间持续而富有成果的交流。Brown 运动由植物学家 Brown 于 1829 年发现，19 世纪众多物理学家对它进行了深入的研究，\(^{(3)}\) 但第一个令人满意的数学模型直到 1905 年才由 Einstein 发明。在粒子沿直线运动的简单情形中，Einstein 的基本假设可以表述如下：若 \(x(t)\) 是粒子在时刻 t 的横坐标，且 \(t_{0} < t_{1} < \cdots < t_{n-1} < t_{n}\)，则相继的位移 \(x(t_{i}) - x(t_{i-1})\)（对 \(1 \leqslant i \leqslant n\)）是独立的 Gauss 随机变量。这里不是详细评述促使 Einstein 理论产生的 J. Perrin 的重要实验工作的地方；对我们的目的而言，只需记住 Perrin 的一个注记，按照这个注记，Brown 运动轨迹的观察使他不可抗拒地联想到“数学家们的没有导数的函数”。这个注记成为 Wiener 的最初火花。

另一个思想潮流发源于气体动理论，它由 Boltzmann 与 Gibbs 在 1870 年至 1900 年间发展。考虑由 N 个质量为 m 的分子组成的气体，处于（绝对）温度 T，并用 \(v_{1},\ldots,v_{N}\) 表示该气体 N 个分子的速度；系统的动能等于

\[
\frac {m}{2} (\mathbf {v} _ {1} ^ {2} + \dots + \mathbf {v} _ {\mathrm{N}} ^ {2}) = 3 \mathrm{N} k \mathrm{T},\tag{1}
\]

其中 k 是 Boltzmann 常数。按照 Gibbs 的想法，分子之间大量的碰撞不容许精确地确定分子的速度，于是方便的做法是在由方程 (1) 定义的 3N 维空间的球面 S 上引入一个概率律 P。“微正则”假设在于假定 P 是球面 S 上在旋转下不变的质量为 1 的测度。此外，Maxwell 速度律指出，分子速度的一个分量的概率律是一个方差为 \(2k\mathrm{T} / m\) 的 Gauss 测度（§6, No.~5, Remark 3）。Borel 似乎是最早在 1914 年观察到：当分子数很大时，Maxwell 律是 Gibbs 假设与球面性质的一个推论。他考虑高维 Euclid 空间中的一个球面 S 以及 S 上在旋转下不变的质量为 1 的测度 P；利用基于 Stirling 公式的经典近似方法，他证明了 P 在一条坐标轴上的投影近似地是 Gauss 的。这些结果稍后由 Gâteaux 与 Lévy (IX, \(a\) ) 加以精确化。给定一个整数 \(m \geqslant 1\) 与一个数 \(r > 0\)，用 \(\mathrm{S}_{m,r}\) 表示形如 \((x_{1}, \ldots, x_{m}, 0, 0, \ldots)\)、满足 \(x_{1}^{2} + \cdots + x_{m}^{2} = r^{2}\) 的序列所成的集合；再用 \(\sigma_{m,r}\) 表示 \(\mathrm{S}_{m,r}\) 上在旋转下不变的质量为 1 的测度。用现代的语言陈述，Gâteaux 与 Lévy 的结果如下：测度序列 \(\sigma_{m,1}\) 紧收敛到原点 \((0, 0, \ldots)\) 处的一个单位质量，而测度序列 \(\sigma_{m,\sqrt{m}}\) 紧收敛到一个形如下式的测度 \(\Gamma\)

\[
d \Gamma (x _ {1}, x _ {2}, \dots) = \prod_ {n = 1} ^ {\infty} d \gamma (x _ {n})
\]

（γ 是 R 上方差为 1 的 Gauss 测度）。

前面的测度 \(\Gamma\) 起着无穷维 Gauss 测度的作用。Lévy 似乎确实曾模糊地希望以内在的方式在每个无穷维 Hilbert 空间上定义一个 Gauss 测度。事实上，正如 Lévy 与 Wiener 所指出的，测度 \(\Gamma\) 在某种意义下 \(^{(4)}\) 在 \(l^{2}\) 的自同构之下是不变的；不幸的是，集合 \(l^{2}\)（由平方可和序列 \((x_{1}, x_{2}, \ldots, x_{n}, \ldots)\) 组成）对 \(\Gamma\) 的测度为零。现在人们知道，只能满足于无穷维 Hilbert 空间上的一个 Gauss 预测度。\(^{(5)}\)

本质的进展要归功于 Wiener：既然在无穷维 Hilbert 空间上没有合理的 Gauss 测度，人们就可以借助取原函数的运算，从一个 Gauss 预测度出发，在连续函数空间上构造一个测度 \(w\)（细节 cf.~§6, No.~7, Th. 1）。我们将扼要地解释 Wiener 对 \(w\) 的原始构造 (X)；它直接受到 Gâteaux 与 Lévy 的关系式 \(\Gamma = \lim_{m \to \infty} \sigma_{m,\sqrt{m}}\) 的影响。对每个整数 \(m \geqslant 1\)，用 \(H_m\) 表示 \(T = ]0,1]\) 上在每个区间 \(\left[ \frac{k-1}{m}, \frac{k}{m} \right]\)（对 \(k = 1,2,\ldots,m\)）上取常值的函数所成的集合，用 \(\pi_m\) 表示 \(R^m\) 中半径为 1 的 Euclid 球面上在旋转下不变的质量为 1 的测度。设 \(f_m\) 是 \(H_m\) 到 \(R^m\) 上的同构，它把取值 \(a_k\)（在区间 \(\left[ \frac{k-1}{m}, \frac{k}{m} \right]\) 上）的函数对应到向量 \((a_1, a_2 - a_1, \ldots, a_m - a_{m-1})\)（Wiener 钟爱的“微分空间”一词即由此而来）；用 \(w_m\) 表示 \(H_m\) 上 \(\pi_m\) 在 \(f_m^{-1}\) 下的像测度。Wiener 把所想要的测度 \(w\) 定义为测度 \(w_m\) 的极限。确切地说，用 \(H\) 表示 \(T\) 上赋以一致收敛拓扑的调节函数所成的集合（有 \(H_m \subset H\)，对每个整数 \(m \geqslant 1\)）；对每个有界的一致连续函数 \(F\)（\(H\) 上的），极限 \(A\{F\} = \lim_{m \to \infty} \int_{H_m} F(x) dw_m(x)\) 存在；接着，Wiener 通过对掷硬币游戏之起伏的精细分析得到某些上界估计，并重新利用 Daniell 所强调的紧性论证，证明此时满足应用 Daniell 扩张定理的条件。由此得出存在一个测度 \(w\)（由 \(C(T)\) 支撑），使得 \(A\{F\} = \int_{C(T)} F(x) dw(x)\)。Wiener 随后能够证明测度 \(w\) 对应于 Einstein 的假设，\(^{(6)}\) 而且他的估计使他能够对 Perrin 关于无导数函数的注记赋予精确的意义：满足阶为 \(\frac{1}{2}\) 的 Lipschitz 条件的函数所成的集合对 \(w\) 是可忽略的（然而，对每个 \(a\)（满足 \(0 < a < \frac{1}{2}\)），几乎每个函数都满足阶为 \(a\) 的 Lipschitz 条件）。

如今，Wiener 测度的众多构造已是熟知的。例如，Paley 与 Wiener 使用了随机 Fourier 级数 (XI, Ch. IX)：对每个实

\[
\begin{array}{l} \int_ {\mathcal {C} (\mathrm{T})} f (x (t _ {1}), \dots , x (t _ {n})) d w (x) = \\ (2 \pi) ^ {- n / 2} \prod_ {i = 1} ^ {n} (t _ {i} - t _ {i - 1}) ^ {- 1 / 2} \int \dots \int f (x _ {1}, \dots , x _ {n}) \exp \left(- \frac {1}{2} \sum_ {i = 1} ^ {n} \frac {(x _ {i} - x _ {i - 1}) ^ {2}}{t _ {i} - t _ {i - 1}}\right) d x _ {1} \dots d x _ {n}, \end{array}
\]

序列 \(\mathbf{a} = (a_n)_{n\geqslant 1}\) 与每个整数 \(m\geqslant 0\)，用下式定义 ]0,1] 上的函数 \(f_{m,\mathbf{a}}\)：

\[
f _ {m, \mathbf {a}} (t) = a _ {1} t + 2 \sum_ {k = 2} ^ {2 ^ {m + 1}} \frac {1}{\pi k} a _ {k - 1} \sin \pi k t;
\]

可以证明，对 \(\Gamma\)-几乎每个序列 \(\mathbf{a}\)，函数序列 \(f_{m,\mathbf{a}}\) 趋于一个连续函数 \(f_{\mathbf{a}}\)，且 \(w\) 是 \(\Gamma\) 在（几乎处处定义的）映射 \(\mathbf{a} \mapsto f_{\mathbf{a}}\) 下的像。后来，Lévy 在 (IX, \(b\) ), \(c\) ) 中给出了一个与我们在 §6, No.~7 中所阐述的非常接近的构造。最后，Kac、Donsker 与 Erdős 大约在 1950 年指明，如何用更一般的测度代替 Wiener 原始构造中的球面测度 \(\pi_m\)（在 \(\mathbf{R}^m\) 上）。他们的结果在 Wiener 测度与概率论的极限定理之间建立了牢固的联系；这些结果后来由 Prokhorov (XIII) 加以完成和系统化，我们将在后面回到这一工作。

这里不是分析由 Wiener 的发现所产生的众多而重要的概率论工作的地方；如今，Brown 运动只表现为 Markoff 过程最重要的一些例子之一。我们只提一下 Kac 把 Wiener 测度应用于求解某些抛物型偏微分方程；这是对量子场论中 Feynmann 思想的一个改编---这又是数学与物理问题相互影响的一个例子。

\specsection{测度的逆极限}

这一理论主要是为满足概率论的需要而发展起来的。关于随机变量有限序列 \(X_{1},\ldots,X_{n}\) 的问题，原则上在知道该序列的律 \(P_{X}\) 时就解决了：这是 \(R^{n}\) 上质量为 1 的一个正测度，使得同时得到不等式 \(a_{1}\leqslant X_{1}\leqslant b_{1},\ldots,a_{n}\leqslant X_{n}\leqslant b_{n}\) 的概率等于 \(\mathrm{P}_{\mathrm{X}}(\mathrm{C})\)，其中 C 是闭方体 \([a_{1},b_{1}]\times\cdots\times[a_{n},b_{n}]\)（在 \(R^{n}\) 中）。在实践中，测度 \(P_{X}\) 或者有离散支撑，或者关于 Lebesgue 测度有密度。当处理随机变量的无穷序列 \((\mathrm{X}_{n})_{n\geqslant1}\) 时，律 \(P_{n}\)（部分序列 \((\mathrm{X}_{1},\ldots,\mathrm{X}_{n})\) 的）对每个整数 \(n\geqslant1\) 都是已知的；这些数据满足一个相容性条件，它表示序列 \((\mathrm{P}_{n})_{n\geqslant1}\) 是测度的一个逆系统（或“射影系统”）。直到大约 1920 年，涉及无穷序列的事件的概率还是以或多或少隐含的方式，由有限情形概率的“自然”极限过渡来定义的；例如，人们假设一局赌博终止的概率是“它至多在 n 步内终止”的概率当 n 趋于无穷时的极限。自然，这样的理论不太协调，而且没有什么能排除“悖论”的出现：同一个概率，按照人们用两个（各自都同样“自然的”）程序中的这一个或那一个来计算，会得到两个不同的估计。

Steinhaus (V) 似乎是最早感到有必要（对掷硬币游戏）不仅考虑逆系统 \((\mathbf{P}_{n})_{n\geqslant1}\)、而且考虑其极限的人。稍早一些，在 1919 年，Daniell (VI, b)) 已经一般地证明了这种逆极限的存在，\(^{(7)}\) 但这一结果在欧洲似乎一直不为人所知。它在 1933 年被 Kolmogoroff 在著作 (XII) 中重新发现，作者在该著作中阐述了概率论的公理化概念。Daniell 与 Kolmogoroff 的证明利用了一个紧性论证，它实质上与我们在 §4, No.~3 的 Th. 2 中所用的相同，并且基于 Dini 定理。

Daniell--Kolmogoroff 定理对随机序列 \((\mathbf{X}_{n})_{n\geqslant1}\) 的情形已无可挑剔，但从 1935 年起由 Kolmogoroff、Feller 与 Doob 开展的随机函数的研究则隐含着完全另一种层次的困难。例如考虑 R 的一个区间 T，表示一个“随机过程”的观察时刻所成的集合；可能轨迹所成的集合是乘积空间 \(R^{T}\)，它被视为部分乘积 \(R^{H}\) 的逆极限，其中 H 取遍 T 的有限子集所成的集合；一般假设给定了测度的一个逆系统 \((\mu_{\mathrm{H}})\)（cf.~§4, No.~2）。Kolmogoroff 定理确实给出了 \(R^{T}\) 上的一个测度，但它定义在一个比 Borel σ-集环小得多的 σ-集环上。\(^{(8)}\) Kolmogoroff 构造的一个变体（它给出拓扑空间上的一个测度）应归功于 Kakutani（Proc. Imp. Acad. Tokyo 19 (1943), 184--188），此后被多次重新发现：把 \(\mu_{H}\) 视为 \(\overline{R}^{H}\) 上由 \(R^{H}\) 支撑的测度；\(^{(9)}\) 紧空间 \(E = \overline{R}^{T}\) 是有限乘积 \(\overline{R}^{H}\) 的逆极限，于是可以把 E 上的测度 \(\mu\) 定义为诸 \(\mu_{H}\) 的逆极限（cf.~Ch. III, §4, No.~5）。然而，这一做法有一个严重的不便；\(\overline{R}^{T}\) 的元素不具备任何能使过程的概率研究取得进展的正则性质---甚至不能简单地消去寄生值 \(\pm\infty\)（由 R 的紧化 \(\overline{R}\) 所引入的）。补救的办法是把测度 \(\mu\)（\(\overline{R}^{T}\) 上的）诱导到一个特定的子空间上（例如 Brown 运动情形中的 \(\mathcal{C}(T)\)）；基本的困难来自这样一个事实：一个函数空间，即使是通常类型的，也不一定是 \(\mu\) -可测的（在 \(\overline{R}^{T}\) 中），甚至连函数空间的选择都可能成问题。\(^{(10)}\)

决定性的一步是 Prokhorov 于 1956 年在一部对随机过程理论产生了决定性影响的著作 (XIII) 中迈出的。通过把 Wiener 在上面分析的文章中所用的方法放进一个适当的公理化形式，他建立了函数空间上测度逆极限的一个一般存在性定理，它是对应于波兰空间的 §4, No.~2 的 Th. 1 的一个特殊情形。

更受限的一类逆系统由 Bochner (XIV) 于 1947 年引入；这是由有限维实向量空间与满射线性映射组成的逆系统。这样一个系统的逆极限可以以自然的方式等同于代数对偶 \(\mathbf{E}^*\)（实向量空间 \(\mathbf{E}\) 的），赋以弱拓扑 \(\sigma (\mathbf{E}^{*},\mathbf{E})\)；相应的测度逆系统有一个极限，即测度 \(\mu\)（定义在一个比 \(\mathbf{E}^*\) 的 Borel σ-集环小得多的 σ-集环上）。Bochner 通过 Fourier 变换完全刻画了这种“预测度”，Fourier 变换是 \(\mathbf{E}\) 上的一个函数。但在 \(\mathbf{E}\) 上没有拓扑时，这一结果几乎无法使用，这时就必须考察把 \(\mu\) 视为拓扑对偶 \(\mathbf{E}'\)（\(\mathbf{E}\) 的）上的测度的可能性。以独立的方式，R. Fortet 与 E. Mourier 在寻求把概率论的某些经典结果（大数定律、中心极限定理）推广到取值于 Banach 空间的随机变量时，也揭示了 Fourier 变换在这类问题中所起的基本作用。但实质性的进展直到 1956 年才到来：Gelfand (XV, \(b\) ) 提出，Fourier 变换的自然框架不是 Banach 空间或 Hilbert 空间，而是核 Fréchet 空间。他猜想这样的空间上每个正型的连续函数是其对偶上某个测度的 Fourier 变换，这一结果不久由 Minlos (XVI) 证明。它的重要性首先在于它适用于广义函数空间，而且几乎全部的函数空间都是广义函数空间的 Borel 子集（因而它构成一个比 \(\mathbf{R}^{\mathrm{T}}\) 好得多的容器）。\(^{(11)}\) 随机广义函数理论是一个蓬勃发展的领域，我们只需请读者参考 Gelfand 与 Vilenkin 的书 (XVII)。

刚才提到的关于逆极限的结果利用了基空间上拓扑的存在性。人们可能会问，在“抽象”测度的情形是否存在类似的理论。Von Neumann 于 1935 年证明了乘积测度在所有情形下的存在性，但 Jessen 与 Andersen (XVIII) 发现的一个反例使每个测度逆系统都容许极限的希望破灭了。人们发现了两种权宜之计：1949 年，C. Ionescu-Tulcea 借助适当分解的存在性建立了可数逆极限的存在，\(^{(12)}\) 这是研究 Markoff 过程的一个非常有趣的结果；此外，人们已经观察到，空间的拓扑只是通过紧子集所成的集合这一中介介入。因此，很自然地试图在抽象理论内部、借助集合的紧类概念把这种情形公理化。这项工作于 1953 年由 Marczewski（他用这种方法建立了一个抽象逆极限定理）与 Ryll-Nardzewski（他处理了测度的分解）完成。\(^{(13)}\)

\specsection{一般拓扑空间上的测度与紧收敛}

对拓扑与测度论之间联系的研究，首先被视为对测度正则性性质、特别是“外”正则与“内”正则的研究；\(^{(14)}\) 在无穷处可数的局部紧空间上，内正则等价于外正则。Lebesgue 对直线上子集测度的构造突出了这两种正则性，而波兰空间上测度的外正则性质似乎在 1935 年前后已为公众所知。但直到 1940 年，A. D. Alexandroff 才在一篇传播因战争而推迟的文章 (XIX) 中突出了内正则的作用，并证明波兰空间上的测度具有这一性质；这一结果后来被 Prokhorov (XIII) 重新发现，并且常被错误地归功于这位作者。直到最近人们才认识到这一性质可以推广到 Souslin 空间；因此，这些空间的重要性大大增加，尤其是因为人们认识到它们的理论可以在没有任何可度量化假设的情况下展开，而且几乎全部的函数空间都是 Souslin 的（多数情形甚至是 Lusin 的）。\(^{(15)}\) 这些就是促使我们在本章强调内正则测度的理由。

测度收敛方式（淡收敛或紧收敛）的定义，最方便的做法是把测度空间与一个连续函数空间做成对偶。A. A. Markoff 推广了 F. Riesz 的一个老结果，于 1938 年建立了 \(\mathcal{C}(\mathrm{X})\) 上正泛函与紧空间 X 上正则测度之间的一一对应。在前面已引用的著作 (XIX) 中，A. D. Alexandroff 把这些结果推广到完全正则空间的情形：他在空间 \(\mathcal{C}^b (\mathrm{X})\)（完全正则空间 \(\mathrm{X},^{(16)}\) 上有界连续函数的）上的正线性形式所成的集合中引入了一个层级，他定义了有界测度的紧收敛，并证明了下面两个定理：

\begin{enumerate}
\def\labelenumi{\alph{enumi})}
\item
  若 X 是波兰空间，则对应于测度的 \(\mathcal{C}^{b}(X)\) 上线性形式所成的集合对序列的弱收敛是闭的；
\item
  若一列有界测度有紧极限，则“没有质量逃逸到无穷远”（这是关于紧收敛的 Prokhorov 定理之逆的一个弱形式）。
\end{enumerate}

从这些丰富的概念与定理中，Prokhorov 提取了对随机过程理论重要的结果，并以简单而醒目的形式呈现它们。在他 1956 年已引用的大著作 (XIII) 中，很大一部分专门讨论波兰空间上的有界正测度；他推广了 Lévy 的一个构造，在质量为 1 的正测度所成的集合上定义了一个距离，使它成为一个波兰空间，然后建立了一个重要的紧性判别法（cf.~§5, No.~5, Th. 1）。与 Prokhorov 相独立，Le Cam (XX) 得到了若干关于测度紧收敛的紧性结果；他对所考虑的空间不作任何可度量化假设，而且在局部紧情形他的结果归结为 Dieudonné 的较早的定理。

\specsection{参考文献}

\begin{enumerate}
\def\labelenumi{(\Roman{enumi})}
\item
  Ch. DE LA VALLÉE POUSSIN, Intégrales de Lebesgue, Fonctions d'ensembles, Classes de Baire, Paris (Gauthier--Villars), 1st. edn., 1916; 2nd. edn., 1936.
\item
  M. FRÉCHET: a) Sur l'intégrale d'une fonctionelle étendue à un ensemble abstrait, Bull. Soc. Math. France, 43 (1915), pp.~248--265; b) Les familles et fonctions additives d'ensembles abstraits, Fund. Math., 4 (1923), pp.~329--365; ibid. 5 (1924), pp.~206-251.
\item
  S. SAKS, Theory of the Integral, 2nd. edn., New York (Stechert), 1937.
\item
  E. BOREL, Les probabilités dénombrables et leurs applications arithmétiques, Rend. Circ. Math. Palermo, 27 (1909), pp.~247-271.
\item
  H. STEINHAUS, Les probabilités dénombrables et leur rapport à la théorie de la mesure, Fund. Math., 4 (1923), pp.~286--310.
\item
  P. J. DANIELL: a) Integrals in an infinite number of dimensions, Ann. of Math., 20 (1918--19), pp.~281--288; b) Functions of limited variation in an infinite number of dimensions, ibid. 21 (1919--20), pp.~30--38.
\item
  B. JESSEN, The theory of integration in a space of an infinite number of dimensions, Acta Math., 63 (1934), pp.~249-323.
\item
  A. EINSTEIN, Investigations on the Theory of the Brownian Movement, New York (Dover), 1956.
\item
  P. LÉVY: a) Leçons d'Analyse Fonctionnelle, Paris (Gauthier-Villars), 1922（第二版于 1951 年由同一出版社以 Problèmes concrets d'Analyse Fonctionnelle 为题出版）; b) Processus stochastiques et mouvement brownien, Paris (Gauthier-Villars), 1948; c) Le mouvement brownien, Mémor. Sci. Math., 126 (1954).
\item
  N. WIENER, Differential space, J. Math. Phys. MIT, 2 (1923), pp.~131--174 (= Selecta, pp.~55--98, Cambridge (MIT Press), 1964).
\item
  R. E. A. C. PALEY 与 N. WIENER, Fourier transforms in the complex domain, Amer. Math. Soc. Colloq. Publ. No.~19, New York, 1934.
\item
  A. N. KOLMOGOROFF, Grundbegriffe der Wahrscheinlichkeit-srechnung, Berlin (Springer), 1933.
\item
  Ju. V. PROKHOROV, Convergence of random processes and limit theorems in probability theory, Theor. Probability Appl., 1 (1956), pp.~156--214.
\item
  S. BOCHNER, Harmonic Analysis and the Theory of Probability, Berkeley (University of California Press), 1960.
\item
  I. M. GELFAND: a) Generalized stochastic processes（俄文）, Dokl. Akad. Nauk SSSR, 100 (1955), pp.~853--856; b) Some problems of functional analysis（俄文）, Uspehi Mat. Nauk, 11 (1956), pp.~3--12 (= Amer. Math. Soc. Transl. (2), 16 (1960), pp.~315--324).
\item
  R. A. MINLOS, Generalized random processes and their extension to a measure（俄文）, Trudy Moskov. Math. Obšč., 8 (1959), pp.~497--518 (= Selected translations in mathematical statistics and probability, III (19), pp.~291--313).
\item
  I. M. GELFAND 与 N. Ya. VILENKIN, Generalized Functions, Vol. IV, New York (Academic Press), 1964（英译本）.
\item
  E. SPARRE-ANDERSSEN 与 B. JESSEN, On the introduction of measures in infinite product sets, Dansk. Vid. Selbskab. Mat. Fys. Medd., 25 (1948), No.~4, pp.~1-7.
\item
  A. D. ALEXANDROFF, Additive set functions in abstract spaces. I (Ch. 1), Mat. Sbornik, 8 (1940), pp.~307--348; II (Chs. 2 and 3), ibid., 9 (1941), pp.~563--628; III (Chs. 4 to 6), ibid., 13 (1943), pp.~169--238.
\item
  L. LE CAM, Convergence in distribution of stochastic processes, Univ. Cal. Publ. Statistics, No.~11 (1957), pp.~207--236.
\end{enumerate}

\specialchapter{记号索引}

参考编号依次表示章、节和小节。

第七章：

\(\gamma_{X}(s)\) , \(\gamma(s)\) : VII, 1, 1.

\(\gamma(s)f\) , \(\gamma(s)\mu\)（f 为函数，\(\mu\) 为测度）: VII, 1, 1.

\(d\mu(s^{-1}x)\) : VII, 1, 1.

\(\delta_{X}(s)\) , \(\delta(s)\) , \(\delta(s)f\) , \(\delta(s)\mu\) , \(d\mu(xs)\) : VII, 1, 1.

\(\check{f}\) , \(\check{\mu}\) , \(d\mu(x^{-1})\)（f 为函数，\(\mu\) 为测度）: VII, 1, 1.

\(\Delta_{G}\) , \(\Delta\) : VII, 1, 3.

mod \(_{G}\varphi\) , mod \(\varphi\)（\(\varphi\) 为自同构）: VII, 1, 4.

\(Z_{p}\)（p 为素数）: VII, 1, 6.

\(K^{+}\)（K 为域）: VII, 1, 10.

mod \(_{K}a\) , mod a（a 为局部紧域 K 的元素）: VII, 1, 10.

\(\mathcal{H}^{\chi}(X)\) , \(\mathcal{H}_{+}^{\chi}(X)\) , \(\mathcal{H}^{1}(X)\) , \(f^{\chi}\) , \(f^{1}\)（X 为局部紧空间，其上有局部紧群 H 作用，\(\chi\) 为 H 在 \(R_{+}^{*}\) 中的连续表示 : VII, 2, 1.

\(f^{b}\) : VII, 2, 2.

\(\lambda^{\sharp}\) , \(\frac{\mu}{\beta}\) , \(\mu/\beta\) : VII, 2, 2.

\(m^{\sharp}\)（m 为向量测度）: VII, 2, 2.

\(T_{J}\) , \(T_{1}(n,K)\) , \(T(n,K)\) , \(T(n,K)^{*}\) : VII, 3, 3.

第八章：

\(^{n}_{i=1}\mu_{i}\) , \(*_{\varphi}(\mu_{i})_{1\leqslant i\leqslant n}\) , \(\mu_{1}*\mu_{2}*\cdots*\mu_{n}\) : VIII, 1, 1.

\(\gamma_{\chi}\) : VIII, 2, 3 与 VIII, 2, 4.

\(\gamma_{\chi,p}\) : VIII, 2, 5.

\(U(\mu)\)（U 为局部紧群 G 的一个表示，\(\mu\) 为 G 上的测度）: VIII, 2, 6.

\(\mathcal{M}^{\rho}(G)\)（G 为局部紧群）: VIII, 3, 1.

\(\mu*\beta f\) , \(\mu*f\)（\(\mu\) 为测度，f 为函数）: VIII, 4, 1.

\(\mathcal{L}(G)\)（G 为局部紧群）: VIII, 4, 5.

\(\mathcal{U}_{s}^{\infty}(G)\)（G 为局部紧群）: VIII, 4, Exer. 21.

第九章：

\(\mathcal{F}_{+}(\mathrm{T}), \mathcal{F}_{+}, f_{\mathrm{A}}, f^{0}\): 预备约定。

\(\pi(p), p_{\mathrm{A}}\) 或 \(p|_{\mathrm{A}}: \mathrm{IX}, 1, 1\).

\(\mathcal{P}(\mathrm{T}; \mathbf{C}), \mathcal{P}(\mathrm{T}; \mathbf{R}), \mathcal{P}(\mathrm{T}), \mathcal{P}_{+}(\mathrm{T}): \mathrm{IX}, 1, 2\).

\(w^{\bullet}(f), \int^{\bullet} f dw, \int^{\bullet} f(t) dw(t): \mathrm{IX}, 1, 2\).

\(w^{\bullet}, w_{\mathrm{K}}^{\bullet}: \mathrm{IX}, 1, 2\).

\(w^{+}, w^{-}, |w|: \mathrm{IX}, 1, 2\).

\(\mu(f), \mu(\mathrm{A}): \mathrm{IX}, 1, 5\).

\(\operatorname{Supp}(\mu): \mathrm{IX}, 1, 6\).

\(\sum_{i \in \mathrm{I}} \mu_i: \mathrm{IX}, 1, 7\).

\(\mu^*(f), \mu^*(\mathrm{A}), \int^* f d\mu, \int^* f(t) d\mu(t): \mathrm{IX}, 1, 9\).

\(\mu^*: \mathrm{IX}, 1, 9\).

\(\overline{\mathcal{L}}^p(\mathrm{T}, \mu), \overline{\mathcal{L}}_F^p(\mathrm{T}, \mu), \mathcal{L}^p(\mathrm{T}, \mu), \mathcal{L}_F^p(\mathrm{T}, \mu)\)（对 \(1 \leqslant p \leqslant +\infty\)）: IX, 1, 10.

\(\overline{\mathcal{L}}_F^p(\mu), \overline{\mathcal{L}}_F^p, \overline{\mathcal{L}}^p, \overline{\mathcal{L}}^p(\mu), \mathcal{L}^p(\mu), \mathcal{L}^p: \mathrm{IX}, 1, 10\).

\(\overline{\mathrm{N}}_p(f), \mathrm{N}_p(f), \overline{\mathcal{N}}_F, \mathcal{N}_F: \mathrm{IX}, 1, 10\).

\(\mathrm{L}_F^p(\mu), \mathrm{L}_F^p: \mathrm{IX}, 1, 10\).

\(\int f d\mu, \mu(f), \int f(t) d\mu(t): \mathrm{IX}, 1, 10\).

\(\mu_X^\bullet, \mu_X, \mu|_{\mathrm{X}}: \mathrm{IX}, 2, 1\).

\(f \cdot \mu: \mathrm{IX}, 2, 2\).

\(\pi(\mu): \mathrm{IX}, 2, 3\).

\(\lambda \otimes \mu: \mathrm{IX}, 2, 5\).

\(\Re(\mathrm{T}), \Re(\mathrm{T}): \text{§3 的约定}\).

\(\mathcal{C}^b(\mathrm{T}; \mathrm{F}), \mathcal{C}^b(\mathrm{T}), \mathcal{C}^b, \mathcal{C}_+^b(\mathrm{T}), \mathcal{C}_+^b: \text{§5 的约定}\).

\(\mathcal{M}^b(\mathrm{T}; \mathbf{C}), \mathcal{M}^b(\mathrm{T}), \mathcal{M}^b, \mathcal{M}_+^b(\mathrm{T}), \mathcal{M}_+^b: \text{§5 的约定}\).

\(\mathcal{L}\mu: \mathrm{IX}, 5, 7\).

\(\mathcal{F}(\mathrm{E}): \mathrm{IX}, 6, 1\).

\(p_V, p_{VW}: \mathrm{IX}, 6, 1\).

\(\mathcal{D}(\mathrm{E}): \mathrm{IX}, 6, 1\).

\(\widetilde{\lambda}: \mathrm{IX}, 6, 1\).

\(u(\mu)\)（\(\mu\) 为预测度）: IX, 6, 2.

\(\mathcal{F}\mu\)（\(\mu\) 为预测度或测度）: IX, 6, 3.

\(\Gamma_Q, \gamma_a: \mathrm{IX}, 6, 5\).

\(\gamma_C: \mathrm{IX}, 6, 6\).

\(\operatorname{Tr}(\mathrm{Q}/\mathrm{H}): \mathrm{IX}, \text{附录}, 1\).

\(u^*: \mathrm{IX}, \text{附录}, 2\).
